\documentclass[11pt]{article}
\usepackage[a4paper,margin=28mm]{geometry}
\usepackage{amsmath,amssymb,amsthm,mathtools}
\usepackage[colorlinks=true,linkcolor=blue,citecolor=blue,urlcolor=blue]{hyperref}
\newtheorem{proposition}{Proposition}[section]
\newtheorem{conjecture}[proposition]{Conjecture}
\theoremstyle{definition}
\newtheorem{definition}[proposition]{Definition}
\newtheorem{example}[proposition]{Example}
\newtheorem{remark}[proposition]{Remark}
\newcommand{\ii}{\mathrm i}
\newcommand{\A}{\mathbb A}
\newcommand{\C}{\mathbb C}
\newcommand{\length}{\operatorname{length}}
\newcommand{\Sym}{\operatorname{Sym}}
\newcommand{\ord}{\operatorname{ord}}
\title{Completeness of the Bethe Ansatz\\through Partition Ideals and Weighted Point Sums}
\author{}
\date{}
\begin{document}
\maketitle

\begin{abstract}
We study an algebraic counting approach to completeness for the periodic
spin-$1/2$ XXX chain. A partition-lattice formula defines depermutation
with local intersection multiplicities, and a binomial alternating sum
corrects the singular-pair contributions. Normalized inhomogeneous
families give a local complex proving nonnegativity of the corrected
lengths without spectral input. At a simple fully reduced zero, direct
cluster elimination also identifies this index with the length of an
explicit residual ideal; a weighted local hypothesis gives a jet-space
Koszul model for that length. Finite flatness of the remainder
families also proves the affine degree formula under the partition-ideal
finiteness hypotheses. A spectral resolution further identifies the
corrected sum with the normalized ordered-root pullback of the
polynomial-kernel ideal. Identifying its pointwise coefficients with the
usual physical-root weights still requires root-simplicity statements;
in particular, the uniform exclusion of repeated singular pairs remains
unproved, although several families are treated explicitly.
\end{abstract}

\section{Background: the XXX Bethe ansatz}

Let $L$ be the length of the periodic spin-$1/2$ XXX chain, with state space
$\mathcal H=(\C^2)^{\otimes L}$. One normalization of its Hamiltonian is
\[
 H=\frac14\sum_{\ell=1}^L
 \bigl(\boldsymbol{\sigma}_{\ell}\cdot
 \boldsymbol{\sigma}_{\ell+1}-1\bigr),
 \qquad \boldsymbol{\sigma}_{L+1}=\boldsymbol{\sigma}_1.
\]
The sector with $n$ down spins has dimension $\binom Ln$.
The algebraic Bethe ansatz constructs candidate eigenvectors
\[
 B(\lambda_1)\cdots B(\lambda_n)|0\rangle
\]
from parameters satisfying
\begin{equation}\label{eq:bethe}
 \left(\frac{\lambda_k+\ii/2}{\lambda_k-\ii/2}\right)^L
 =
 \prod_{j\ne k}\frac{\lambda_k-\lambda_j+\ii}{\lambda_k-\lambda_j-\ii},
 \qquad k=1,\ldots,n.
\end{equation}
Throughout these notes we assume $L\ge2n$, where $n$ is the number of
Bethe roots. The number of highest-weight states in this
sector, counted with their vector-space dimensions, is
\begin{equation}\label{eq:target}
 d(L,n)=\binom Ln-\binom L{n-1},
\end{equation}
%Indeed, each irreducible $\mathfrak{sl}_2$ summand contributes one state to each of its weight spaces; taking the difference of two consecutive sector dimensions extracts the number of highest-weight states of the specified weight. This is a target dimension for completeness, rather than a count of distinct energy eigenvalues.

One would expect that the Bethe ansatz reconstructs all highest-weight
vectors. For simplicity, rescale the parameters by
\[
 \lambda_j=\frac{\ii}{2}x_j,
 \qquad x_j=\frac{2\lambda_j}{\ii}.
\]
The singular values $\lambda_j=\pm\ii/2$ become $x_j=\pm1$, and the
interaction shifts $\pm\ii$ become $\pm2$.
After clearing denominators and discarding a nonzero common scalar,
the polynomial is
\begin{equation}\label{eq:polynomial}
	P(x_1,\ldots,x_n)=
	(x_1+1)^L\prod_{j\ne1}(x_1-x_j-2)
	-
	(x_1-1)^L\prod_{j\ne1}(x_1-x_j+2).
\end{equation}
It is symmetric in its last $n-1$ arguments. Set
\[
 P_i(x)=P(x_i,x_1,\ldots,\widehat{x_i},\ldots,x_n).
\]
The periodic Bethe equations are $P_i=0$ for $i=1,\ldots,n$.

First we require distinct coordinates; repeated coordinates are called
\emph{collisions}. The Bethe creation operators commute, so we also
identify configurations differing by a permutation. Solutions containing
$x_1=1$ and $x_2=-1$ are called \emph{singular}. For distinct remaining
roots, physical singular solutions of the periodic system require the
additional regularization condition
\begin{equation}\label{eq:physical}
 \left[-\prod_{j=3}^n
 \frac{x_j+1}{x_j-1}\right]^L=1.
\end{equation}
As an example for $n=2$, the empty product is $1$,
and the condition requires $L$ to be even.
See \cite{NW2013} for regularizations and
\cite{HNS2013} for the completeness counting problem.

\section{The counting strategy}
Assume that the set of ordered solutions of $P_1=\cdots=P_n=0$ is finite,
and write
\[
 X=\{x^1,\ldots,x^s\}\subseteq\C^n
\]
for this set. In the target count, physical solutions have pairwise
distinct coordinates and yield nonzero Bethe vectors, using
regularization at singular solutions. Each physical permutation orbit
therefore has $n!$ elements. Assign the target weights
\[
 w_i=\begin{cases}
 1/n!,&x^i\text{ is a physical solution},\\
 0,&\text{otherwise}.
 \end{cases}
\]
The desired completeness statement is
\begin{equation}\label{eq:target-weights}
 \sum_{i=1}^s w_i=d(L,n).
\end{equation}

We work with formal sums of points with rational coefficients:
\[
 Z=\sum_i c_i[x^i],\qquad c_i\in\mathbb Q.
\]
Here $[x^i]$ is a formal symbol for the point $x^i$, and equality of
such sums means equality of the coefficient at every point. We write
$\deg Z=\sum_i c_i$ for the total weight, allowing negative coefficients.
All operations below are operations on these formal sums.
For a map $\phi$ of point sets, we use the notation
\[
 \phi_*\!\left(\sum_p c_p[p]\right)=\sum_p c_p[\phi(p)],
\]
adding coefficients when several points have the same image. In
particular, passing from ordered tuples to unordered configurations
simply adds the weights over each permutation orbit. We write
$\Sym^r(\A^1)$ for the set of unordered configurations of $r$ complex
numbers, allowing repetitions.

For a zero-dimensional ideal $I\subseteq\C[x_1,\ldots,x_n]$, define
its global and local lengths by
\[
 \begin{aligned}
 \ell(I)&=\dim_\C\bigl(\C[x_1,\ldots,x_n]/I\bigr),\\
 \ell_p(I)&=\dim_\C\bigl(\C[x_1,\ldots,x_n]/I\bigr)_{\mathfrak m_p},
 \end{aligned}
\]
where $\mathfrak m_p$ is the maximal ideal of the point $p$, and
$\ell_p(I)=0$ if $p\notin V(I)$. The associated formal sum is
\[
 Z(I)=\sum_{p\in V(I)}\ell_p(I)[p],\qquad
 \deg Z(I)=\ell(I).
\]
Thus its coefficients are nonnegative integers, and its total weight
is the dimension of a finite-dimensional quotient algebra.

Our strategy is to express the target weights as rational linear
combinations of ordinary local intersection lengths. We seek ideals
$I_j$ and coefficients $m_j\in\mathbb Q$ such that each $I_j$ is
generated by $n$ polynomials, is zero-dimensional, and satisfies
$V(I_j)\subseteq X$. Set $Z_j=Z(I_j)$ and write
\[
 Z_j=\sum_{i=1}^s l_{ij}[x^i],\qquad
 l_{ij}=\ell_{x^i}(I_j),
\]
with $l_{ij}=0$ when $x^i\notin V(I_j)$. We require the formal sum
\begin{equation}\label{eq:weight-cycle}
 Z_{\text{Bethe}}=\sum_jm_jZ_j
\end{equation}
to equal $\sum_iw_i[x^i]$; equivalently,
\begin{equation}\label{eq:local-weight-decomposition}
	\sum_j m_jl_{ij}=w_i\qquad\text{for every }i.
\end{equation}
Taking degrees then gives
\begin{equation}\label{eq:global-weight-decomposition}
 \sum_iw_i=\sum_jm_j\ell(I_j),\qquad
 \ell(I_j)=\deg Z_j=\sum_i l_{ij}.
\end{equation}

The generators of $I_j$ allow us to study their global
intersection numbers by BKK theory or a suitable compactification.

The proposed approach has two steps:
\begin{enumerate}
 \item To remove collision and permutation redundancy we make use of M\"obius inversion on the partition lattice. We call this process {\em depermutation}.
 \item Construct an excess formal sum on the singular-pair locus and subtract it.
 The resulting formal sum should satisfy the pointwise identities
 \eqref{eq:local-weight-decomposition} and hence the degree identity
 \eqref{eq:target-weights}.
\end{enumerate}
The first step supplies explicit formal sums $Z_j$ and rational coefficients
$m_j$ for depermutation. The second modifies these coefficients and adds
formal sums on the singular locus so that the resulting local sums are the
target weights $w_i$. The advantage is the passage from the local
identities \eqref{eq:local-weight-decomposition} to the global calculation
\eqref{eq:global-weight-decomposition}.

We use the following notation for the auxiliary formal sums. Write $F_i=P_i$,
and let $A,B\subseteq[n]$ be disjoint, with either $A=B=\varnothing$ or
both $A$ and $B$ nonempty. Let $\pi=\{S_1,\ldots,S_m\}$ be a partition
of $R=[n]\setminus(A\cup B)$, and choose a representative $r_\alpha\in
S_\alpha$ for each block. For $t\in R$, define
the reduced polynomial by
\begin{equation}\label{eq:reduced-fixed-equation}
 F_t^{\mathrm{red}}(x_R)
 :=\frac{\left.F_t(x)\right|_{x_a=1\ (a\in A),\,
                                      x_b=-1\ (b\in B)}}{(x_t+1)^{|A|}(x_t-1)^{|B|}}.
\end{equation}
This is well-defined: the two terms in the numerator are both divisible
by $(x_t+1)^{|A|}(x_t-1)^{|B|}$, since $L\ge2n$.
For $A=B=\varnothing$, set $F_t^{\mathrm{red}}=F_t$.
The dependence on $A,B$ is suppressed in the notation. Define
\begin{equation}\label{eq:fixed-partition-ideal}
 \begin{aligned}
 I_{A,B;\pi}=\bigl(&x_a-1\ (a\in A),\quad x_b+1\ (b\in B),\\
                 &F_{r_\alpha}^{\mathrm{red}}\ (1\le\alpha\le m),\\
                 &x_t-x_{r_\alpha}\quad
                  (1\le\alpha\le m,\ t\in S_\alpha\setminus\{r_\alpha\})
                  \bigr).
 \end{aligned}
\end{equation}
Here one reduced equation $F_{r_\alpha}^{\mathrm{red}}=0$ is retained for each block, together
with the equalities within that block. The ideal has $|A|+|B|$ fixed-root
equations, $m$ Bethe equations, and $n-|A|-|B|-m$ equality equations,
for a total of $n$ specified generators. Symmetry makes it independent
of the chosen representatives. Let
\[
 Z_{A,B;\pi}=Z(I_{A,B;\pi})
 =\sum_{p\in V(I_{A,B;\pi})}\ell_p(I_{A,B;\pi})[p],
\]
whenever $I_{A,B;\pi}$ is zero-dimensional. Its support is contained in
$X$: the original equations at the fixed roots vanish automatically when
both signs occur, and each remaining original equation is
$(x_t+1)^{|A|}(x_t-1)^{|B|}F_t^{\mathrm{red}}$ after fixing $A,B$. Within each block,
the reduced equations agree after imposing the equalities. We abbreviate
$Z_{\varnothing,\varnothing;\pi}$ to $Z_\pi$.

Our aim is to find rational coefficients $m_{A,B;\pi}$ such that
\[
 Z_{\text{Bethe}}=\sum_{A,B,\pi}m_{A,B;\pi}Z_{A,B;\pi}.
\]
This is the notation $j=(A,B;\pi)$ in
\eqref{eq:local-weight-decomposition}; its local lengths are
$l_{i,(A,B;\pi)}=\ell_{x^i}(I_{A,B;\pi})$.
Throughout the note, $I$ denotes an ideal and $Z$ a formal sum of points.
All auxiliary ideals $I_{A,B;\pi}$ use the reduced equations
\eqref{eq:reduced-fixed-equation}; the associated sums $Z_{A,B;\pi}$
therefore need no additional superscript. The original polynomials
$F_t$ are retained for studying multiplicities in the full Bethe system.

\subsection{Depermutation}
The general principle for removing collisions is M\"obius inversion on the
partition lattice. To illustrate it, let $\Omega$ be a finite set and
$S\subseteq\Omega^n$ be invariant under coordinate permutations. For a
partition $\pi$ of $[n]$, let $S_\pi$ consist of the tuples in $S$ satisfying
$x_i=x_j$ whenever $i,j$ belong to the same block of $\pi$; coordinates in
different blocks are still allowed to coincide. Order the partition lattice
$\Pi_n$ by refinement, with $\hat0$ the partition into singletons. Then
the number of permutation orbits in $S$ with pairwise distinct
coordinates is
\[
 \frac{1}{n!}\sum_{\pi\in\Pi_n}\mu(\hat0,\pi)\,\#S_\pi,
 \qquad
 \mu(\hat0,\pi)=\prod_{B\in\pi}(-1)^{|B|-1}(|B|-1)!.
\]


We now replace these finite sets by formal sums of points recording local lengths.
Consider a general polynomial $f(x_1,\ldots,x_n)$ symmetric in
its last $n-1$ arguments. Set
\[
 F_i(x)=f(x_i,x_1,\ldots,\widehat{x_i},\ldots,x_n).
\]
For a partition $\pi=\{S_1,\ldots,S_m\}$ choose a representative
$r_a\in S_a$ for each block, and set
\[
 I_\pi=\bigl(F_{r_a},\ x_b-x_{r_a}:
       1\le a\le m,\ b\in S_a\setminus\{r_a\}\bigr).
\]
There are $m$ polynomial equations and $n-m$ equality equations, hence
$n$ specified generators. Set $Z_\pi=Z(I_\pi)$.
Symmetry makes the ideal independent of the representatives. Its support
is contained in the full solution set: on each block diagonal all the
corresponding $F_i$ agree. For $f=P$ this support is contained in $X$;
for general $f$ use its full solution set in place of $X$.

Assume these ideals are zero-dimensional. In the notation above, the index $j$
is now the partition $\pi$, and
\[
 l_{i\pi}=\ell_{x^i}(I_\pi),\qquad
 m_\pi=\frac{\mu(\hat0,\pi)}{n!},\qquad
 Z_\pi=\sum_i l_{i\pi}[x^i].
\]

\begin{definition}
\label{def:depermutation}
The {\em depermutation sum} of $f$ is
\begin{equation}\label{eq:cycle}
 Z^{\mathrm{dep}}_f=\frac{1}{n!}
 \sum_{\pi\in\Pi_n}
 \mu(\hat0,\pi)Z_\pi
 =\sum_\pi m_\pi Z_\pi.
\end{equation}
\end{definition}

We now return to the XXX polynomial $f=P$, so that $F_i=P_i$.
To formulate the singular correction, extend depermutation to systems
with fixed singular roots. For $0\le k\le\lfloor n/2\rfloor$,
let $A,B\subseteq[n]$ be disjoint subsets with $|A|=|B|=k$, and put
$R=[n]\setminus(A\cup B)$ and $r=|R|=n-2k$.
For these balanced fixed sets, use the reduced polynomials from
\eqref{eq:reduced-fixed-equation}. Assuming the ideals
$I_{A,B;\pi}$ are zero-dimensional, set
\begin{equation}\label{eq:reduced-fixed-pair-depermutation}
 \begin{aligned}
 Z^{\mathrm{dep}}_{A,B}
 &=\frac1{r!}\sum_{\pi\in\Pi(R)}
       \mu(\hat0,\pi)Z_{A,B;\pi},\\
 Z^{\mathrm{dep}}_k
 &=
 \frac{(k!)^2(n-2k)!}{n!}\sum_{\substack{A,B\subseteq[n],\ A\cap B=\varnothing\\
                   |A|=|B|=k}}Z^{\mathrm{dep}}_{A,B}.
 \end{aligned}
\end{equation}
Here $\Pi(R)$ is the partition lattice of the remaining labels.
Depermutation acts only on the $r$ free coordinates; the sum over $A,B$
restores invariance under all coordinate permutations. The factor
$(k!)^2r!/n!$ averages over the $n!/((k!)^2r!)$ choices of $A,B$.
For $r=0$, the partition lattice has one element, with M\"obius value
one, and $0!=1$. In particular, $Z^{\mathrm{dep}}_0=Z^{\mathrm{dep}}_P$.

\begin{conjecture}[Alternating completeness formula]
\label{conj:alternating-completeness}
Assume $L\ge2n$ and that the reduced partition ideals are
zero-dimensional. The formal sum
\begin{equation}\label{eq:central-alternating-cycle}
 \boxed{
 Z^{\mathrm{cand}}_{L,n}
 :=\sum_{k=0}^{\lfloor n/2\rfloor}(-1)^k\binom Lk
 Z^{\mathrm{dep}}_k
 }
\end{equation}
has the target physical weights, that is,
\[
 Z^{\mathrm{cand}}_{L,n}=\sum_{i=1}^s w_i[x^i].
\]
Moreover,
\[
 \deg Z^{\mathrm{cand}}_{L,n}
 =d(L,n)=\binom Ln-\binom L{n-1}.
\]
\end{conjecture}

The affine degree assertion follows from
Proposition~\ref{prop:affine-completeness-degree} under the stated
finiteness hypotheses. The pointwise identification with the target
physical weights remains the conjectural part.

\subsection{A generalized $TQ$ relation and depermutation}
\label{sec:general-tq}

Fix a polynomial $f$ symmetric in its last $n-1$ arguments, and write
$F_i(x)=f(x_i,x_1,\ldots,\widehat{x_i},\ldots,x_n)$.
Let $e_j$ and $\hat e_j$ denote the elementary symmetric polynomials
in $x_1,\ldots,x_n$ and $x_2,\ldots,x_n$, respectively, and set $e_0=1$.
The map from ordered roots to their elementary symmetric functions is
\[
 q:\A^n\longrightarrow\Sym^n(\A^1)\simeq\A^n,
 \qquad x\longmapsto e(x).
\]
Using $e_1,\ldots,e_n$ as coordinates on the target, define
\[
 Q(t;e)=\sum_{j=0}^n(-1)^je_jt^{n-j},
\]
Write $Z^{\mathrm{dep}}$ for the depermutation sum of this system.

The $TQ$ relation transforms equations in the roots $x$ into equations
in their elementary symmetric functions $e$. We construct a polynomial
$\Phi(t;e)$ whose divisibility by $Q(t;e)$
encodes the equations at distinct roots. By symmetry, there is a
polynomial $g$ such that
\[
 f(x_1,\ldots,x_n)=g(x_1;\hat e_1,\ldots,\hat e_{n-1}).
\]
Set
\begin{align*}
	E_j(t;e)&=\sum_{k=0}^j(-t)^ke_{j-k},\qquad 1\le j<n,\\
	\Phi(t;e)&=g\bigl(t;E_1(t;e),\ldots,E_{n-1}(t;e)\bigr).
\end{align*}

The elementary symmetric identities give
$E_j(x_1;e(x))=\hat e_j$, hence
\[
 \Phi(x_1;e(x))=g(x_1;\hat e_1,\ldots,\hat e_{n-1})=F_1(x).
\]
If the coordinates are distinct, the equations $F_i(x)=0$ therefore
give $\Phi(x_i;e(x))=0$ for every $i$, by symmetry. Thus they are
equivalent to $Q\mid\Phi$ on this locus.

Since $Q$ is monic, polynomial division gives unique polynomials $U,R$
with coefficients in $\C[e_1,\ldots,e_n]$ such that
\begin{equation}\label{eq:general-tq-division}
	\Phi(t;e)=U(t;e)Q(t;e)+R(t;e),\qquad
	R(t;e)=\sum_{j=0}^{n-1}\rho_j(e)t^j.
\end{equation}

Define the \emph{generalized $TQ$ ideal}
\begin{equation}\label{eq:general-tq-ideal}
 J=(\rho_0,\ldots,\rho_{n-1})\subseteq\C[e_1,\ldots,e_n].
\end{equation}

The formal sum associated to the $TQ$ ideal is the pushforward of the
depermutation sum to $\Sym^n(\A^1)$:
\[
 q_*Z^{\mathrm{dep}}=Z(J).
\]
Here we assume the relevant ideals are zero-dimensional. Informally,
one can see this by perturbing all equations coherently as
$F_i\mapsto F_i+h(x_i)$, where $h$ is a generic small polynomial of
degree less than $n$. The local intersection lengths are conserved
under this perturbation. For generic $h$, the nearby solutions are
simple, and M\"obius inversion cancels those with repeated coordinates.
Each remaining unordered configuration corresponds to one solution of
the perturbed $TQ$ ideal and has $n!$ ordered representatives, each of
weight $1/n!$. This proves the identity for the perturbed systems;
specializing back recovers the local multiplicities on both sides.

For the XXX system, we use the following additional input.

\paragraph{Distinct-root assumption away from the singular values.}
Let $L\ge2n$, and suppose a monic polynomial $Q$ of degree $n$ admits
a polynomial $T$ satisfying Baxter's relation
\begin{equation}\label{eq:xxx-tq-relation}
 T(t)Q(t)=(t+1)^LQ(t-2)+(t-1)^LQ(t+2).
\end{equation}
Write
\[
 Q(t)=(t-1)^a(t+1)^bQ_{\mathrm{reg}}(t),
 \qquad Q_{\mathrm{reg}}(1)Q_{\mathrm{reg}}(-1)\ne0.
\]
We assume that $Q_{\mathrm{reg}}$ is squarefree. Equivalently, every
root of $Q$ outside $\{1,-1\}$ is simple; repetitions at the singular
values are allowed.

The exclusion of repeated roots through polynomiality of $T$ is
discussed in \cite[Introduction, eq.~(1.3) and footnote~2]{HJJ2024},
where numerical evidence is reported and a rigorous proof for the
spin-$1/2$ XXX chain is stated to be unavailable to the authors.
We use the precise version above, including configurations with
singular roots, as an assumption rather than a proved result.
Together with the pushforward identity, it rules out collisions away
from $\pm1$ in the support of $Z^{\mathrm{dep}}$. The same conclusion
holds for the reduced balanced fixed-pair systems: if $H$ satisfies
their $TQ$ relation after fixing $k$ pairs, then
$Q=(t^2-1)^kH$ satisfies~\eqref{eq:xxx-tq-relation}.

\subsection{Virtual degrees and the combinatorial identity}

\begin{definition}[Virtual degree assignment]
\label{def:virtual-degree}
For a partition $\pi$ of $[n]\setminus(A\cup B)$ with $m$ blocks, set
\begin{equation}\label{eq:virtual-partition-degree}
 \boxed{\deg_{\mathrm{vir}}Z_{A,B;\pi}=(L+r-2)^m.}
\end{equation}
When $r=0$, assign degree one.
\end{definition}

The assignment has the following BKK interpretation. Put
$r=n-|A|-|B|$ and $D=L+r-2$. For $r\ge1$, retain the ambient monomial
supports of the reduced equations when identifying coordinates within
each block, before combining coefficients that may cancel. Let
$P_\alpha\subseteq\mathbb R_{\ge0}^m$ be the downward closure of the
convex hull of the retained exponents for block $\alpha$. The ambient
equation has total degree at most $D$ and a nonzero pure term
$x_{r_\alpha}^D$. Consequently,
\[
 [0,De_\alpha]\subseteq P_\alpha\subseteq D\Delta_m,
 \qquad
 \Delta_m=\{u\in\mathbb R_{\ge0}^m:\textstyle\sum_\alpha u_\alpha\le1\}.
\]
Normalize mixed volume by
$\operatorname{MV}(\Delta_m,\ldots,\Delta_m)=1$. Monotonicity gives
\[
 \begin{aligned}
 D^m
 &=\operatorname{MV}([0,De_1],\ldots,[0,De_m])\\
 &\le\operatorname{MV}(P_1,\ldots,P_m)
 \le\operatorname{MV}(D\Delta_m,\ldots,D\Delta_m)=D^m.
 \end{aligned}
\]
Thus the mixed volume equals $D^m$. BKK identifies this with the generic
solution count when all monomials at lattice points of the enlarged
polytopes are allowed independent coefficients; generically the
solutions have no zero coordinates, so this is also the affine count
\cite[Chapter~3]{Sturmfels}.

This mixed-volume count agrees with the B\'ezout number $D^m$ for
$m$ hypersurfaces of degree $D$ in $\mathbb P^m$. We may therefore use
the standard compactification $\mathbb A^m\subset\mathbb P^m$ to
study the virtual count. For a block equation $G_\alpha(y)$, use its
degree-$D$ homogenization
\[
 \widehat G_\alpha(Y_1,\ldots,Y_m,T)
 =T^D G_\alpha(Y_1/T,\ldots,Y_m/T),
\]
retaining degree $D$ even when cancellation lowers the actual affine
degree. The hyperplane $T=0$ then records the contribution at infinity.
If the homogenized equations have a finite common zero set, their
total length is $D^m$, including local lengths at infinity. If they
have a positive-dimensional common set at infinity, $D^m$ remains the
virtual count obtained by a generic degree-$D$ deformation, rather than
a sum of isolated local lengths for the special equations.

\begin{proposition}[Virtual completeness count]
\label{prop:virtual-completeness}
For $|A|=|B|=k$ and $r=n-2k$,
\begin{equation}\label{eq:virtual-depermutation-degree}
 \deg_{\mathrm{vir}}Z^{\mathrm{dep}}_{A,B}
 =\deg_{\mathrm{vir}}Z^{\mathrm{dep}}_k
 =\binom{L+r-2}{r}.
\end{equation}
Consequently,
\begin{equation}\label{eq:virtual-completeness-degree}
 \boxed{
 \deg_{\mathrm{vir}}Z^{\mathrm{cand}}_{L,n}
 =\sum_{k=0}^{\lfloor n/2\rfloor}(-1)^k\binom Lk
       \binom{L+n-2k-2}{n-2k}
 =\binom Ln-\binom L{n-1}.}
\end{equation}
\end{proposition}

\begin{proof}
The partition-lattice identity
\[
 \sum_{\pi\in\Pi_r}\mu(\hat0,\pi)D^{|\pi|}
 =D(D-1)\cdots(D-r+1)
\]
follows by counting injective maps from $[r]$ to a set of $D$ elements:
there are $D^{|\pi|}$ maps constant on each block of $\pi$, and
M\"obius inversion removes all coincidences. Dividing by $r!$ gives
\[
 \deg_{\mathrm{vir}}Z^{\mathrm{dep}}_{A,B}
 =\frac1{r!}\sum_{\pi\in\Pi_r}\mu(\hat0,\pi)D^{|\pi|}
 =\binom Dr=\binom{L+r-2}{r}.
\]
For $r=0$, both sides are one. The average defining
$Z^{\mathrm{dep}}_k$ has $n!/((k!)^2r!)$ equal-degree terms, so its
virtual degree is the same.

For the final identity, use formal power series:
\[
 \begin{aligned}
 \left(\sum_{k=0}^L(-1)^k\binom Lk t^{2k}\right)
 \left(\sum_{r\ge0}\binom{L+r-2}{r}t^r\right)
 &=\frac{(1-t^2)^L}{(1-t)^{L-1}}\\
 &=(1-t)(1+t)^L.
 \end{aligned}
\]
Taking the coefficient of $t^n$ proves
\eqref{eq:virtual-completeness-degree}, with
$\binom L{-1}=0$ if $n=0$.
\end{proof}

Thus the virtual assignment gives the desired number. To prove
Conjecture~\ref{conj:alternating-completeness}, we must connect this
virtual count with the affine solutions and establish the prescribed
weight at every point.

\subsection{Plan for a proof of Conjecture~\ref{conj:alternating-completeness}}\label{section:plan}
For $p\in X$, define its exact singular sets
\[
 A(p)=\{j:p_j=1\},\qquad B(p)=\{j:p_j=-1\},
\]
and let $\pi(p)$ be the equality partition of the remaining labels:
two labels belong to the same block precisely when their coordinates
are equal. Write $a=|A(p)|$ and $b=|B(p)|$. If one of $a,b$ is
positive, the original Bethe equation at that fixed root forces the
other to be positive as well. Hence either $a=b=0$ or $a,b\ge1$.
These exact labels classify the affine configurations. Write
$c_k(p)=\operatorname{coeff}_p Z^{\mathrm{dep}}_k$.

Here list some claims that's needed to prove the conjecture:
\begin{enumerate}
	\item Solution at infinity: there are some cancellation such that $P$ do not have desired degree when you make a series of substitution like $x_i=x_j$. As an example, when $L=6,n=3$, $P(x,x,x)=48x^5+160x^3+48x$ but the expected degree is $7$. We claimed those "solution at infinity" will be cancelled at depermutation.
	\item Every $Z^{dep}_{k}$ contains only points with singleton $\pi$. 
	\item Every $Z^{dep}_{k}$ contains only points with $a=b$.
	\item \textbf{Exactly one singular pair: $a=b=1$ and $\pi=\hat 0$.}
	Only $k=0,1$ can contribute, so prove
	\begin{equation}\label{eq:singular_weight}
	c_0(p)-Lc_1(p)=w(p)
	\end{equation}
	where $w(p)$ is the target weight above.
	Writing $e_p$ for the local length of the reduced fixed-pair
	system, the precise local-length claim is
	\begin{equation}\label{eq:single-pair-length-criterion}
	 \ell_p=
	 \begin{cases}
	 Le_p+1,&p\text{ satisfies the physical condition }\eqref{eq:physical},\\
	 Le_p,&\text{otherwise}.
	 \end{cases}
	\end{equation}
	For a simple reduced solution, $e_p=1$.

	\item \textbf{Multiple balanced singular pairs: $a=b\ge2$ and $\pi=\hat 0$.}
	These points have repeated coordinates and hence target weight zero.
	Prove
	\[
	\sum_{k=0}^{a}(-1)^k\binom Lk c_k(p)=0,
	\]

\end{enumerate}

\paragraph{Reduction using the distinct-root assumption.}
The distinct-root assumption substantially reduces these tasks.
Together with the $TQ$ interpretation of depermutation, it excludes
every affine configuration with a collision outside $\pm1$, for the
full system and for each reduced balanced fixed-pair system. Thus only
configurations with $\pi(p)=\hat0$ need further consideration.

At infinity, the expected cancellation has a similar interpretation:
several roots tend to the same point $\infty\in\mathbb P^1$.
Indeed, in the unpartitioned reduced system with $r\ge1$ free roots,
the coefficient of the pure term $x_i^{L+r-2}$ in its own equation is
$2(L-2n+2)\ne0$. Hence its projective equations have no boundary
solution with just one nonzero root coordinate. Every boundary
solution therefore has at least two coordinates representing roots
at infinity. If the distinct-root exclusion and the
$TQ$/depermutation correspondence extend to a compatible
compactification, these repeated infinite roots have zero depermuted
weight as well. This extension remains a boundary claim to justify:
the affine distinct-root assumption alone does not establish
cancellation at infinity.

For the remaining affine configurations, the $TQ$ relation also
excludes unequal singular multiplicities $a\ne b$, by comparing
vanishing orders at $1$ and $-1$. The same argument applies after
fixing any $k$ balanced pairs, so $c_k(p)=0$ for every
$0\le k\le\min(a,b)$ at an unbalanced configuration.
Consequently, among singular affine configurations, only
the balanced cases $a=b$ remain: the physical-weight identity for
$a=b=1$ and the alternating cancellation for $a=b\ge2$.

\section{Small numbers of Bethe roots}
Before addressing the general claims, we examine small values of $n$.
These calculations distinguish the global degree identities from the
local multiplicity statements and the remaining physical interpretation.
Throughout, we retain the assumption $L\ge2n$.

\subsection{Two roots: $n=2$}

As explained in the section \ref{section:plan},  only interesting thing to show is that on $(1,-1)$ the equality \ref{eq:singular_weight} holds, we start from deciding the multiplicity in $Z_0$.
Put $x=1+u$ and $y=-1+v$. The first equation gives
\[
 v=u-2^{2-L}u^L+O(u^{L+1}).
\]
Substituting into the second gives
\[
 F_2\bigl(1+u,-1+v(u)\bigr)
 =4\bigl((-1)^L-1\bigr)u^L+O(u^{L+1}).
\]
If $L$ is odd the leading coefficient is $-8$ and the ordinary local
multiplicity is $L$. If $L$ is even, the next term is
\[
 F_2\bigl(1+u,-1+v(u)\bigr)
 =-4L\,u^{L+1}+O(u^{L+2}),
\]
so the multiplicity is $L+1$. Neither singular point lies on the
diagonal, and its coefficient in $Z^{\mathrm{dep}}_1$ is $1/2$.
Therefore, at either $p=(1,-1)$ or $p=(-1,1)$,
\[
 \operatorname{coeff}_p Z^{\mathrm{cand}}_{L,2}
 =\frac{\ell_p(I_{\mathrm{full}})-L}{2}
 =\begin{cases}
 0,&L\text{ odd},\\
 1/2,&L\text{ even}.
 \end{cases}
\]
The total remaining weight of the unordered singular pair is thus $0$
for odd $L$ and $1$ for even $L$, exactly matching the periodic singular
admissibility condition \eqref{eq:physical}; see also \cite{NW2013}.

\subsection{Three roots: $n=3$}
\label{sec:n3-infinity}
Let $L\ge6$. Only $k=0,1$ occur in the alternating formula.
The polynomial is
\[
 \begin{aligned}
 P(x,y,z)={}&(x+1)^L(x-y-2)(x-z-2)\\
 &-(x-1)^L(x-y+2)(x-z+2),
 \end{aligned}
\]
with $F_1=P(x,y,z)$, $F_2=P(y,x,z)$, and $F_3=P(z,x,y)$.
Proposition~\ref{prop:virtual-completeness} already gives its virtual
degree, so we focus on the boundary contribution and the singular
local weights.

\paragraph{Cancellation at infinity.}
As explained we can compactify by $\mathbb{P}^3$, Homogenize the three original equations to their ambient degree
$L+1$ before imposing partition equalities. The highest-degree part is
\[
 F_1^{[L+1]}=2x^{L-1}q_x,\qquad
 q_x=(L-4)x^2-(L-2)x(y+z)+Lyz,
\]
with the other expressions obtained by permutation. The only common
point at infinity is $P=[1:1:1:0]$. Indeed, a zero coordinate is
incompatible with the other two equations, and
\[
 q_x-q_y=(x-y)\bigl((L-4)(x+y)-(2L-2)z\bigr)
\]
and its cyclic analogues force at least two coordinates to coincide.
Substitution then forces all three to coincide. For a two-block
partition, say $x=y$, the first leading form restricts to
$4x^L(z-x)$; its common point at infinity with the remaining equation
is again $P$.

Write $\widehat F_i$ for the degree-$(L+1)$ homogenizations.
Permutation symmetry gives
\[
 \widehat F_1-\widehat F_2=(X-Y)K_{12},\qquad
 \widehat F_1-\widehat F_3=(X-Z)K_{13},\qquad
 K_{12}(P)=K_{13}(P)=-12.
\]
Thus, in the local ring at $P$, the full system forces $X=Y=Z$.
The same holds for each two-block system. On the diagonal, in the
chart $X=1$,
\[
 \widehat F_1(1,1,1,T)
 =4T\bigl((1+T)^L-(1-T)^L\bigr)
 =8LT^2+O(T^4).
\]
Hence the local length at infinity is two for every partition type
$(1,1,1)$, $(2,1)$, and $(3)$. Their total contribution to
$Z^{\mathrm{dep}}_0$ is
\[
 \frac16\bigl(2-3\cdot2+2\cdot2\bigr)=0.
\]
This proves the required cancellation for every $L\ge6$.

\paragraph{Unbalanced singular collisions.}
At $p=(1,-1,1)$, we will prove that the depermutation kills the weights. Symmetry gives
\[
 F_1-F_3=(x_1-x_3)K,\qquad K(p)=-2^{L+1}\ne0.
\]
Consequently the full local ideal equals the partition ideal
$(F_1,F_2,x_1-x_3)$. Their equal local lengths cancel with M\"obius
coefficients $1,-1$, so $c_0(p)=0$. Also $c_1(p)=0$, because the
remaining coordinate would be a root $1$ of $Q_L$. Permutation symmetry
and simultaneous sign reversal give the same conclusion at every
configuration with two $1$'s and one $-1$, or vice versa. Computing their
common ordinary local length is unnecessary for this cancellation.

\paragraph{The transverse singular multiplicity.}
Let $p=(1,-1,z)$ be a solution with $z\ne\pm1$. We verify the
local-length criterion~\eqref{eq:single-pair-length-criterion} from
Section~\ref{section:plan}.
The reduced equation is
\[
 Q_L(z)=0,\qquad
 Q_L(w)=(w+1)^{L-1}(w-3)-(w-1)^{L-1}(w+3),
\]
since $F_3(1,-1,w)=(w^2-1)Q_L(w)$. Here the reduced local length is
$e=\ord_zQ_L$, and $\ell_p$ denotes the full local length.

Set $x_1=1+u$, $x_2=-1+v$, and $x_3=w$. Eliminating $v$ by $F_1=0$
gives
\[
 v=u+2^{2-L}\frac{3-w}{1+w}u^L+O(u^{L+1}).
\]
Write the remaining equations as $F_2=u^LH(u,w)$ and $F_3=G(u,w)$.
Their restrictions to $u=0$ are
\[
 \begin{aligned}
 H(0,w)&=\frac4{w+1}
 \bigl((w+1)(w+3)-(-1)^L(w-1)(w-3)\bigr),\\
 G(0,w)&=(w^2-1)Q_L(w).
 \end{aligned}
\]
At a root of $Q_L$, the condition $H(0,z)=0$ is equivalent to
\[
 \left[-\frac{z+1}{z-1}\right]^L=1,
\]
which is the physical singular condition~\eqref{eq:physical}.
Indeed, the reduced equation gives
$((z+1)/(z-1))^{L-1}=(z+3)/(z-3)$, and combining these two
relations yields exactly the numerator of $H(0,z)$.

If $p$ is nonphysical, $H$ is a unit in the local ring. The remaining
ideal is therefore $(u^L,G)$, of length $\ell_p=Le$, as required.

If $p$ is physical, $H(0,z)=0$ forces $z=0$ for even $L$, or
$z^2=-3$ for odd $L$. At these roots $e=1$. The first-order
elimination of $G=0$ gives $w(u)=z+\alpha u+O(u^2)$, and substitution
in $H$ gives $H(u,w(u))=h_1u+O(u^2)$, with
\[
 \begin{array}{c|c|c}
 &\alpha&h_1\\ \hline
 L\text{ even},\ z=0
 &-\dfrac4{3L-4}&\dfrac{36L(L-4)}{3L-4}\\[6pt]
 L\text{ odd},\ z^2=-3
 &-\dfrac2{L-2}&\dfrac{4L(L-3)(z+3)}{L-2}
 \end{array}
\]
In both cases $h_1\ne0$ for $L\ge6$, so the full local multiplicity
is $L+1$. Therefore
\[
 \boxed{
 c_0(p)-Lc_1(p)
 =\begin{cases}
 1/6,&p\text{ is physical},\\
 0,&p\text{ is nonphysical}.
 \end{cases}}
\]
Thus the local-length criterion, and hence the desired pointwise
identity, holds at every distinct-coordinate singular solution for $n=3$.

\subsection{Four roots: the double singular pair $n=4$}
\label{sec:n4-double-pair}

Let $n=4$, $L\ge8$, and $p_0=(1,1,-1,-1)$, and write
$Z^{\mathrm{dep}}_{L,4}=Z^{\mathrm{dep}}_0$. We prove that the corrected
weight at this point is zero. By the multiplicity of the
unordered configuration we mean the sum of the coefficients of
$Z^{\mathrm{dep}}_{L,4}$ over the six distinct permutations of $p_0$.
This multiplicity is before any singular-excess subtraction.

\begin{proposition}\label{prop:n4-double-pair}
The unordered multiplicity of $\{1,1,-1,-1\}$ after depermutation is
\begin{equation}\label{eq:n4-double-pair}
 \boxed{M_L=
 \begin{cases}
 \displaystyle\binom L2,&L\text{ even},\\[2pt]
 \displaystyle\binom{L+1}2,&L\text{ odd}.
 \end{cases}}
\end{equation}
The coefficient at a single ordered point $p_0$ in the normalized formal sum
$Z^{\mathrm{dep}}_{L,4}$ is $M_L/6$.
\end{proposition}

\begin{proof}
We compute the depermutation multiplicity in symmetric coordinates.
This also avoids computing four large ordinary local intersection
lengths separately.

\paragraph{Depermutation and the polynomial remainder equations.}
For a monic quartic $Q$, let $r_Q$ be the remainder modulo $Q$ of
\[
 R_Q(T)=(T+1)^LQ(T-2)+(T-1)^LQ(T+2).
\]
The four coefficients of $r_Q$ define a system in the four coefficients
of $Q$. We will show below that it is isolated at
$Q_0(T)=(T^2-1)^2$, and compute its local length $M_L$. The full and partition local lengths are taken under our standing
finiteness assumption.

This local length is the unordered depermutation multiplicity. To see
this, use the common polynomial perturbation $h_a$ from
Section~\ref{sec:general-tq}, now of degree at most three.
The remainder system becomes $r_Q+h_a=0$. In a finite local family, a
generic small $h_a$ gives $M_L$ simple quartics near $Q_0$, each with
distinct roots. Indeed, this system is the fiber of the map
$Q\mapsto-r_Q$; generic fibers are reduced, and the discriminant
hypersurface of quartics has lower-dimensional image in a finite local
neighborhood. Each such quartic has two roots near $1$ and two near
$-1$, and gives $2!\,2!=4$ ordered configurations near $p_0$.
All other nearby solutions have collisions and cancel in the M\"obius
sum, as in Section~\ref{sec:general-tq}.
Thus the coefficient at $p_0$ is $4M_L/24=M_L/6$, and its orbit has
total coefficient $M_L$.

\paragraph{Separating the two clusters.}
Use the local factorization
\[
 Q(T)=A(T-1)B(T+1),\qquad
 A(u)=u^2-au+b,\qquad B(v)=v^2-cv+d.
\]
The coefficients $(a,b,c,d)$ are local coordinates near $Q_0$: the two
factors of $Q_0$ are coprime, so multiplication is locally invertible.
Divisibility by $Q$ is equivalent to divisibility at each of these
coprime quadratic factors. After division by units, the two congruences
are
\begin{equation}\label{eq:n4-cluster-congruences}
 \begin{aligned}
 B(u)+u^LU(u)&\equiv0\pmod{A(u)},\\
 A(v)+v^LV(v)&\equiv0\pmod{B(v)},
 \end{aligned}
\end{equation}
where
\[
 \begin{aligned}
 U(u)&=\frac{A(u+2)B(u+4)}{(u+2)^LA(u-2)},\\
 V(v)&=\frac{A(v-4)B(v-2)}{(v-2)^LB(v+2)}.
 \end{aligned}
\]
The denominators are units in the corresponding completed quadratic
algebras. The two coefficients of the first congruence solve uniquely
for $c,d$ as formal power series in $a,b$, by the implicit function
theorem.

Give $a,u,v$ weight one and $b$ weight two. The remainder of $u^L$
modulo $A$ has its linear coefficient of weight $L-1$ and its constant
coefficient of weight $L$. Consequently the solution of the first
congruence satisfies
\[
 c=a+O_{\mathrm w}(L-1),\qquad d=b+O_{\mathrm w}(L),
\]
where $O_{\mathrm w}(j)$ denotes terms of weight at least $j$.
In particular, for the low-weight calculation that follows we may put
$B=A$ inside $U,V$.

\paragraph{The leading equations and parity.}
Put $C=2^{4-L}$. Direct expansion to weight one gives
\[
 \begin{aligned}
 U(u)&=C\left(1+\frac{5-L}{2}u-\frac54a+O_{\mathrm w}(2)\right),\\
 V(u)&=(-1)^LC\left(1+\frac{L-5}{2}u+\frac54a
                    +O_{\mathrm w}(2)\right).
 \end{aligned}
\]
After eliminating $c,d$, the leading part of the second congruence is
obtained by taking the remainder modulo $A$ of $u^L(U(u)+V(u))$.
The corrections from $B-A$ have higher weight and do not affect these
initial equations, since $L\ge8$.

If $L$ is even, the constant terms add. Up to a nonzero scalar, the
two initial equations are the coefficients of
\[
 \operatorname{rem}_A u^L.
\]
Their weights are $L-1$ and $L$. They have only the common zero
$(a,b)=(0,0)$, since $A\mid u^L$ forces both roots of $A$ to be zero.

If $L$ is odd, the constant terms cancel. Up to a nonzero scalar, the
two initial equations are the coefficients of
\[
 \operatorname{rem}_A\left((5-L)u^{L+1}-\frac52a u^L\right).
\]
Their weights are $L$ and $L+1$. These also have only the common zero
$(0,0)$. Indeed, if $A$ had two distinct nonzero roots, both would have
to be the same root of the linear factor
$(5-L)u-\tfrac52a$, which is impossible. If its roots were $0$ and
$a\ne0$, evaluation at $u=a$ would give
$(\tfrac52-L)a^{L+1}\ne0$. If it had a nonzero double root $r$, then
$a=2r$ and evaluation at $u=r$ would give $-Lr^{L+1}\ne0$.

In both cases the weighted initial equations form a zero-dimensional
regular sequence, so higher-weight terms do not change the local length.
For weights $1,2$ and equation weights $d_1,d_2$, that length is
$d_1d_2/2$: the Hilbert series is
\[
 \frac{(1-z^{d_1})(1-z^{d_2})}{(1-z)(1-z^2)},
\]
whose value at $z=1$ is $d_1d_2/2$. This proves isolation of the remainder
system and yields $L(L-1)/2$ for even $L$, and $L(L+1)/2$ for odd $L$,
as claimed.

\end{proof}


\paragraph{The reduced one-pair contribution at the double pair.}
To finish the weight calculation at $p_0$, fix one positive and one
negative root. The two remaining reduced equations are $G(x,y)=0$ and
$G(y,x)=0$, where
\[
 G(x,y)=(x+1)^{L-1}(x-3)(x-y-2)
       -(x-1)^{L-1}(x+3)(x-y+2).
\]
Let $b_L$ be their local length at $(1,-1)$. Put $x=1+u$, $y=-1+v$,
and $\eta=u-v$. The first equation gives
\[
 \eta=\frac{4C(u)}{A(u)-C(u)}
 =-2^{4-L}u^{L-1}
   \left(1+\frac{5-2L}{4}u+O(u^2)\right),
\]
where $A(u)=(2+u)^{L-1}(u-2)$ and $C(u)=u^{L-1}(u+4)$.
Substitution into the second equation yields
\[
 G(y,x)=
 \begin{cases}
 32u^{L-1}+O(u^L),&L\text{ even},\\
 8(2L-5)u^L+O(u^{L+1}),&L\text{ odd}.
 \end{cases}
\]
Hence $b_L=L-1$ for even $L$ and $b_L=L$ for odd $L$.

There are four choices of the fixed positive and negative labels at
$p_0$. Each contributes $b_L/2$ before the averaging factor $2!/4!$,
so $c_1(p_0)=b_L/6$. Fixing both pairs gives $c_2(p_0)=1/6$, while
Proposition~\ref{prop:n4-double-pair} gives $c_0(p_0)=M_L/6$.
Consequently,
\[
 \boxed{
 \operatorname{coeff}_{p_0}Z^{\mathrm{cand}}_{L,4}
 =\frac16\left(M_L-Lb_L+\binom L2\right)=0.}
\]
For even $L$ the numerator is
$\binom L2-L(L-1)+\binom L2=0$; for odd $L$ it is
$\binom{L+1}2-L^2+\binom L2=0$.
Thus the double singular pair has the required zero weight for both
parities, as do all its permutations.

\section{Root at infinity: general case}
\label{sec:general-infinity}

Use homogeneous coordinates $[X_1:\cdots:X_n:W]$ on $\mathbb P^n$,
where $W=x_{n+1}$ and $x_i=X_i/W$ on the affine chart. Set
$D=L+n-2$, and let $\widehat F_i$ be the degree-$D$ homogenization
of $F_i$. The cancellation of the terms of degree $D+1$ gives
\[
 \widehat F_i(X,W)=\frac1W\left(
 (X_i+W)^L\prod_{j\ne i}(X_i-X_j-2W)
 -(X_i-W)^L\prod_{j\ne i}(X_i-X_j+2W)\right).
\]
The numerator is divisible by $W$, so this is a polynomial. On the
hyperplane at infinity, its restriction is
\begin{equation}\label{eq:leading-equations-infinity}
 \widehat F_i(X,0)=2X_i^{L-1}\left(
 L\prod_{j\ne i}(X_i-X_j)
 -2X_i\sum_{j\ne i}\prod_{h\ne i,j}(X_i-X_h)
 \right).
\end{equation}

\begin{proposition}[Collisions at infinity]
\label{prop:collisions-infinity}
Assume $L\ge2n$. If $[X_1:\cdots:X_n:0]$ is a common zero of
$\widehat F_1,\ldots,\widehat F_n$, then there are distinct indices
$i,j$ such that $X_i=X_j\ne0$. In particular, every solution at
infinity lies on a collision diagonal.
\end{proposition}

\begin{proof}
Suppose instead that all nonzero coordinates are pairwise distinct.
Let $S=\{i:X_i\ne0\}$, write $s=|S|\ge1$, and let $z=n-s$ be
the number of zero coordinates. For $i\in S$, every factor
$X_i-X_j$ is nonzero. Dividing
\eqref{eq:leading-equations-infinity} by
$2X_i^{L-1}\prod_{j\ne i}(X_i-X_j)$ therefore gives
\[
 L=2\sum_{j\ne i}\frac{X_i}{X_i-X_j}
   =2z+2\sum_{\substack{j\in S\\j\ne i}}
                     \frac{X_i}{X_i-X_j}.
\]
Sum over $i\in S$. Each unordered pair in $S$ contributes one to
the double sum, since
\[
 \frac{X_i}{X_i-X_j}+\frac{X_j}{X_j-X_i}=1.
\]
Consequently,
\[
 Ls=2zs+s(s-1),\qquad
 L=2z+s-1=n+z-1\le2n-2.
\]
This contradicts $L\ge2n$.
\end{proof}

The same argument applies after fixing $k$ balanced singular pairs.
There are $r=n-2k$ free roots, whose reduced equations are
\[
 G_i(x)=(x_i+1)^{L-k}(x_i-3)^k
                \prod_{j\ne i}(x_i-x_j-2)
       -(x_i-1)^{L-k}(x_i+3)^k
                \prod_{j\ne i}(x_i-x_j+2),
\]
where the products run over the other free roots. Their degree is
$L+r-2$, and their leading forms are
\eqref{eq:leading-equations-infinity} with $n$ replaced by $r$ and
the coefficient $L$ inside the parentheses replaced by $L-4k$.
Since $L-4k\ge2r$, the same contradiction shows that every boundary
solution has two equal nonzero free-root coordinates. For $r=1$
there are no boundary solutions, and $r=0$ has no free-root boundary.

Thus the collision at infinity can be proved directly, without the
affine distinct-root assumption. The equality $X_i=X_j\ne0$ means
that the two roots have the same nonzero leading coordinate as they
diverge. This support statement does not yet prove cancellation of
their local multiplicities in the depermutation sum; that requires
comparing the full and partition intersections along this boundary.
\section{One singular pair: $a=b=1$}
\label{sec:one-singular-pair}

Let $n\ge2$, $L\ge2n$, and
$p=(1,-1,z_3,\ldots,z_n)$ have pairwise distinct coordinates.
We retain the standing finiteness assumptions. Put
\[
 I_{\mathrm{full}}=(F_1,\ldots,F_n),\qquad
 I_{\mathrm{red}}=I_{\{1\},\{2\};\hat0},\qquad
 e_p=\ell_p(I_{\mathrm{red}}).
\]
We prove the weight identity without assuming that $e_p=1$.
The physical input is the standard nonzero singular Bethe-vector
regularization of \cite[Section~3]{NW2013}, together with simultaneous
diagonalizability of the homogeneous XXX transfer matrices
\cite[Theorem~12.6.1]{MTV2009}. Only diagonalizability is used, not the
count of physical solutions or simplicity of the joint spectrum.

\begin{proposition}\label{prop:one-pair-weight}
At $p$, the full local length is
\begin{equation}\label{eq:one-pair-general-length}
 \ell_p(I_{\mathrm{full}})=Le_p+\chi_{\mathrm{phys}}(p),
\end{equation}
where $\chi_{\mathrm{phys}}(p)$ is one if
\eqref{eq:physical} holds and zero otherwise. Consequently,
\[
 \boxed{c_0(p)-Lc_1(p)=\frac{\chi_{\mathrm{phys}}(p)}{n!}=w(p).}
\]
\end{proposition}

\begin{proof}
We separate the singular factor and then compute the residual length.

\paragraph{Separating the singular factor.}
Write $x_1=1+u$, $x_2=-1+v$, and $x_j=y_j$ for $j\ge3$.
Set
\[
 C(u,y)=\prod_{j=3}^n(u-1-y_j),\qquad
 D(u,y)=\prod_{j=3}^n(u+3-y_j).
\]
Since $C(0,z)\ne0$, the first equation solves for $v$ exactly as
\[
 v=u+u^Ld(u,y),\qquad
 d(u,y)=-\frac{4D(u,y)}{(2+u)^LC(u,y)-u^LD(u,y)}.
\]
After this substitution, write
\[
 F_2=u^LH(u,y),\qquad F_j=G_j(u,y)\quad(3\le j\le n),
\]
with $H,G_j$ regular near $(0,z)$. Work in the completed local ring
\[
 S=\C[[u,y_3-z_3,\ldots,y_n-z_n]],\qquad
 \mathcal R=S/(G_3,\ldots,G_n).
\]
The full and reduced local algebras
are respectively $\mathcal R/(u^LH)$ and $\mathcal R/(u)$: at $u=0$,
$G_j=(y_j^2-1)F_j^{\mathrm{red}}$, with $y_j^2-1$ a unit.
They have finite length, so $\mathcal R$ is a one-dimensional complete
intersection and $u$ is a non-zero-divisor. The exact sequence
\[
 0\longrightarrow\mathcal R/(H)
 \xrightarrow{\ \times u^L\ }\mathcal R/(u^LH)
 \longrightarrow\mathcal R/(u^L)\longrightarrow0
\]
therefore gives
\begin{equation}\label{eq:one-pair-length-splitting}
 \ell_p(I_{\mathrm{full}})
 =Le_p+\dim_\C\mathcal A,\qquad
 \mathcal A:=S/(H,G_3,\ldots,G_n).
\end{equation}
Here $\mathcal A=0$ when $H$ is a unit.

\paragraph{The physical compatibility condition.}
Put $r=n-2$. Direct substitution gives
\[
 H(0,y)=4(-1)^r\left[
 -\prod_{j=3}^n(y_j+3)
 +(-1)^L\prod_{j=3}^n
             \frac{(y_j-1)(y_j-3)}{y_j+1}\right].
\]
There are no other exact strings: a finite chain of distinct roots
separated by $2$ must end at $-1$ and $1$, hence consists of this pair
alone; see \cite[Lemma~1]{GJ2020}. In particular, $z_j\ne\pm3$
and $z_i-z_j\ne\pm2$ for $i,j\ge3$. Thus all factors in the
following ratios are nonzero. Multiplying the reduced equations
cancels the interactions between the remaining roots and yields
\[
 \left(\prod_{j=3}^n\frac{z_j+1}{z_j-1}\right)^{L-1}
 =\prod_{j=3}^n\frac{z_j+3}{z_j-3}.
\]
Combining this with the expression for $H(0,z)$ gives
\begin{equation}\label{eq:one-pair-compatibility}
 H(0,z)=0
 \quad\Longleftrightarrow\quad
 \left[-\prod_{j=3}^n\frac{z_j+1}{z_j-1}\right]^L=1.
\end{equation}
At a nonphysical point, $H$ is consequently a unit, and
\eqref{eq:one-pair-length-splitting} proves the assertion.

\paragraph{The residual length at a physical point.}
Suppose now that \eqref{eq:physical} holds. The algebra $\mathcal A$
is finite and local, with residue field $\C$. Denote by
$\mathcal B(t)$ and $\mathsf t(t)$ the polynomial creation operator
and transfer matrix in the rescaled normalization. The vector
\[
 \Psi(u,y)=u^{-L}\mathcal B(1+u)\mathcal B(-1+v(u,y))
                    \prod_{j=3}^n\mathcal B(y_j)|0\rangle
\]
is regular at $u=0$, since $v-u=O(u^L)$. Its value at $(0,z)$ is
the nonzero regularized physical Bethe vector; see
\cite[Section~3]{NW2013} and \cite[Lemmas~3--4]{GJ2020}.

Let
\[
 Q_{\mathcal A}(t)=(t-1-u)(t+1-v)\prod_{j=3}^n(t-y_j)
 \in\mathcal A[t].
\]
All original equations vanish in $\mathcal A$, and the roots remain
distinct modulo its maximal ideal. Thus the $TQ$ quotient is a
polynomial $\tau_{\mathcal A}(t)\in\mathcal A[t]$:
\[
 \tau_{\mathcal A}(t)Q_{\mathcal A}(t)
 =(t+1)^LQ_{\mathcal A}(t-2)+(t-1)^LQ_{\mathcal A}(t+2).
\]
The off-shell transfer-matrix identity also gives, over $\mathcal A$,
\begin{equation}\label{eq:one-pair-artin-eigenvector}
 \mathsf t(t)\Psi=\tau_{\mathcal A}(t)\Psi.
\end{equation}
Indeed, divide that identity by $u^L$ before passing to the quotient.
The unwanted term at the first root is zero by $F_1=0$; that at the
second root contains $F_2/u^L=H$. Each other unwanted term contains
$G_j$ and a Bethe vector still containing the singular pair, which is
divisible by $u^L$. Hence all unwanted terms vanish in $\mathcal A$.

The transfer matrices admit a common eigenbasis, by
\cite[Theorem~12.6.1]{MTV2009} and the $\mathfrak{sl}_2$
decomposition of the spin space. In this basis, some component of
$\Psi$ has nonzero residue and is therefore a unit of $\mathcal A$.
Equation~\eqref{eq:one-pair-artin-eigenvector} then forces
$\tau_{\mathcal A}(t)=\tau_p(t)$, the fixed complex eigenvalue
polynomial at $p$.

For this fixed $\tau_p$, the monic degree-$n$ polynomial solution of
the $TQ$ relation is unique, even over $\mathcal A$. To see this, put
\[
 C_m=L(L-1)-4m(L-m+1).
\]
The degree-$n$ solution gives
$\tau_p(t)=2t^L+C_nt^{L-2}+O(t^{L-3})$. If a nonzero difference
of two monic solutions has degree $m<n$ and leading coefficient $b$,
then the coefficient of $t^{L+m-2}$ in its $TQ$ equation is
\[
 (C_m-C_n)b=4(n-m)(L+1-n-m)b\ne0,
\]
because $L\ge2n$. This is impossible, so
$Q_{\mathcal A}=Q_p$. Since the roots of $Q_p$ are distinct,
evaluating at each corresponding root shows that
$u=0$ and $y_j=z_j$ in $\mathcal A$. Thus $\mathcal A=\C$ and
its length is one. This proves~\eqref{eq:one-pair-general-length}.

\paragraph{Normalization of the weights.}
At a distinct-coordinate point, only the discrete partition
contributes to depermutation. Moreover, there is exactly one choice
of the fixed positive and negative labels. Consequently,
\[
 c_0(p)=\frac{\ell_p(I_{\mathrm{full}})}{n!},\qquad
 c_1(p)=\frac{e_p}{n!},\qquad c_k(p)=0\quad(k\ge2).
\]
Substitution gives the claimed physical weight. In particular, when
$e_p=1$, the ordinary local length is $L+1$ at a physical point and
$L$ otherwise, recovering the small-$n$ calculations.
\end{proof}

\subsection{A short proof at a simple reduced zero}
\label{sec:simple-reduced-one-pair}

When the reduced zero is simple, the local-length
criterion~\eqref{eq:single-pair-length-criterion} has a shorter
proof. Only the elimination and compatibility calculations above are
needed; the Artin-algebra argument can be replaced by a first-order
calculation. We retain pairwise distinct coordinates and $L\ge2n$.

\begin{proposition}\label{prop:simple-reduced-one-pair}
Suppose $p=(1,-1,z_3,\ldots,z_n)$ is a simple zero of
$I_{\mathrm{red}}$. Then
\[
 \ell_p(I_{\mathrm{full}})=L+\epsilon,\qquad
 \epsilon:=\dim_\C S/(H,G_3,\ldots,G_n)\ge0,
\]
where $S,H,G_j$ are obtained by the elimination above. The number
$\epsilon$ is the local intersection multiplicity of $H=0$ with
the smooth curve $G_3=\cdots=G_n=0$, and is positive if and only if
the physical condition~\eqref{eq:physical} holds. At a physical
point this intersection is simple, so $\epsilon=1$ and
\eqref{eq:single-pair-length-criterion} holds with $e_p=1$.
Only the last assertion uses the standard nonzero regularized
singular Bethe vector; no theorem of MTV is needed.
\end{proposition}

\begin{proof}
Set $x_1=1+u$. As above, $F_1=0$ solves for
$x_2=-1+v(u,y)$, with $v=u+O(u^L)$. The Jacobian of
$(G_3,\ldots,G_n)$ in the $y$ variables is invertible at $(0,z)$:
its rows are the reduced Jacobian rows multiplied by the nonzero
factors $z_j^2-1$. The implicit function theorem therefore gives a
unique local curve $y=y(u)$ solving these equations. The full local
algebra becomes
\[
 \C[[u]]/(u^Lh(u)),\qquad h(u)=H(u,y(u)).
\]
Consequently,
\begin{equation}\label{eq:simple-pair-residual-multiplicity}
 \epsilon=\dim_\C\C[[u]]/(h(u))
          =\ord_{u=0}h(u),\qquad
 \ell_p(I_{\mathrm{full}})=L+\epsilon.
\end{equation}
The standing finiteness assumption ensures that $h$ is not identically
zero. We assign vanishing order zero to a unit, so the formula also
covers $h(0)\ne0$. This identifies $\epsilon$ as an ordinary zero
multiplicity on the smooth local curve; it is independent of the
choice of local parameter or multiplication of $H$ by a unit.
Since $y(0)=z$, equation~\eqref{eq:one-pair-compatibility} now gives
\[
 \boxed{\epsilon>0
 \quad\Longleftrightarrow\quad h(0)=H(0,z)=0
 \quad\Longleftrightarrow\quad
 \left[-\prod_{j=3}^n\frac{z_j+1}{z_j-1}\right]^L=1.}
\]
Thus positivity and the equivalence with the physical condition follow
from local elimination alone.

At a physical point it remains to prove $h'(0)\ne0$. Suppose instead
that $h(u)=O(u^2)$. Along the curve let $Q(u;t)$ be the monic root
polynomial, and let $\tau(u;t)$ be the polynomial quotient in its
$TQ$ division. The standard regularization gives a holomorphic vector
\[
 \Psi(u)=u^{-L}\mathcal B(1+u)\mathcal B(-1+v(u,y(u)))
                  \prod_{j=3}^n\mathcal B(y_j(u))|0\rangle,
 \qquad \Psi(0)\ne0;
\]
see \cite[Section~3]{NW2013} and \cite[Lemma~3]{GJ2020}.
The off-shell identity~\cite[eqs.~(7)--(8)]{GJ2020}, divided by
$u^L$, yields
\[
 \bigl(\mathsf t(t)-\tau(u;t)\bigr)\Psi(u)=O(u^2).
\]
Indeed, all unwanted terms vanish except the one containing
$F_2/u^L=h(u)$. Replacing the rational $TQ$ expression by its
polynomial quotient changes it only by $O(u^{L+2})$.

The transfer matrices have a common orthogonal eigenbasis by an
elementary Hermiticity argument. In the original spectral parameter,
use
\[
 \begin{aligned}
 \mathsf L_{a\ell}(\lambda)
 &=\lambda I+\frac{\ii}{2}\sum_{\alpha=1}^3
                         \sigma_a^\alpha\sigma_\ell^\alpha,\\
 \mathsf L_{a\ell}(\lambda)^\dagger
 &=\sigma_a^2\mathsf L_{a\ell}(\lambda)^{t_a}\sigma_a^2
                  \qquad(\lambda\in\mathbb R).
 \end{aligned}
\]
Here $t_a$ is transpose in the auxiliary space. Entries at distinct
sites commute, so reversing the product and taking the auxiliary
trace gives $\mathsf t(\lambda)^\dagger=\mathsf t(\lambda)$ for
real $\lambda$. Its commuting polynomial coefficients are therefore
simultaneously diagonalizable; the spectral rescaling preserves this
eigenbasis.

Differentiate the off-shell identity at $u=0$:
\[
 (\mathsf t(t)-\tau(0;t))\Psi'(0)
       =\dot\tau(t)\Psi(0),\qquad
 \dot\tau(t)=\left.\partial_u\tau(u;t)\right|_{u=0}.
\]
Projection onto the joint eigenspace containing $\Psi(0)$ gives
$\dot\tau(t)=0$. Differentiating the $TQ$ relation then shows that
$\dot Q(t)=\partial_uQ(0;t)$ solves the same homogeneous $TQ$
equation as $Q(0;t)$, but has degree less than $n$. This is
impossible unless $\dot Q=0$: for a nonzero leading coefficient
$b$ in degree $m<n$, the highest uncancelled coefficient is
\[
 4(n-m)(L+1-n-m)b\ne0.
\]
However, direct differentiation at $t=1$ gives
\[
 \dot Q(1)=-2\prod_{j=3}^n(1-z_j)\ne0,
\]
a contradiction. Thus $h'(0)\ne0$, and the length at a physical
point is $L+1$, proving the claim.
\end{proof}

\paragraph{The physical polynomial and its local multiplicity.}
Write
\[
 A(y)=\prod_{j=3}^n(y_j+1),\qquad
 B(y)=\prod_{j=3}^n(y_j-1),\qquad
 \mathcal P(y)=(-A(y))^L-B(y)^L.
\]
Thus $\mathcal P(z)=0$ is precisely the physical condition. The
minus sign is outside the product: putting it inside every factor
would instead give $(-1)^{(n-2)L}A(y)^L-B(y)^L$, which need not
define the same condition.

There is an exact multiplicity interpretation using a natural lift
of $\mathcal P$ to the local curve above. Set $v(u)=v(u,y(u))$ and
\[
 \kappa(u)=\frac{2+u}{v(u)-2}\frac{v(u)}u,\qquad
 \widetilde{\mathcal P}(u)
 =\bigl(\kappa(u)A(y(u))\bigr)^L-B(y(u))^L.
\]
Since $v(u)=u+O(u^L)$, the quotient $v(u)/u$ is regular and
$\kappa(0)=-1$. In particular,
$\widetilde{\mathcal P}(0)=\mathcal P(z)$. In fact,
\begin{equation}\label{eq:epsilon-physical-polynomial}
 \boxed{\epsilon=\ord_{u=0}\widetilde{\mathcal P}(u).}
\end{equation}
To check this, put $\eta=u-v(u)$ and
\[
 C_2(u)=\prod_{j=3}^n\bigl(v(u)-3-y_j(u)\bigr).
\]
Multiply the equations for the remaining roots to cancel their mutual
interaction factors, and use $F_1=0$ to eliminate $\eta$. The
cancelled second equation then gives the identity on the curve
\[
 h(u)=-\frac{(4+\eta)C_2(u)(v(u)-2)^L}
                  {(2+u)^L A(y(u))^L}\,
                  \widetilde{\mathcal P}(u).
\]
The prefactor is a unit, since the remaining roots avoid $\pm1$
and $\pm3$. This proves~\eqref{eq:epsilon-physical-polynomial}
directly, without spectral input.

The lift matters: one cannot in general replace
$\widetilde{\mathcal P}(u)$ by $\mathcal P(y(u))$. For $n=2$
and even $L$, the latter is identically zero, whereas the former
has a simple zero and $\epsilon=1$. On the fixed reduced scheme,
the proved simplicity statement does give the numerical identity
\[
 \epsilon=\dim_\C
 \bigl(\mathcal O_{V(I_{\mathrm{red}}),p}/(\mathcal P)\bigr),
\]
because this local ring is $\C$. This last description records
whether the reduced point satisfies the physical condition; the
lifted polynomial is what computes the residual intersection
multiplicity before simplicity is proved.

\section{Multiple balanced singular pairs}
\label{sec:multiple-balanced-pairs}

Let $p$ have exactly $a$ coordinates equal to $1$, $a$ equal to $-1$,
and $r=n-2a$ other, pairwise distinct coordinates $z_1,\ldots,z_r$.
Write
\[
 H(t)=\prod_{j=1}^r(t-z_j),\qquad
 Q_k(t)=(t^2-1)^{a-k}H(t)\quad(0\le k\le a).
\]
The reduced $TQ$ numerator after fixing $k$ pairs is
\[
 \Phi_k(t;Q)=(t+1)^{L-k}(t-3)^kQ(t-2)
             +(t-1)^{L-k}(t+3)^kQ(t+2).
\]
Let $J_k$ be the ideal of its remainder coefficients modulo a variable
monic polynomial $Q$ of degree $n-2k$, and let $\nu_k(p)$ be its local
length at $Q=Q_k$.

\paragraph{Why the balanced points satisfy $TQ$.}
The original polynomial equations do not in general imply divisibility
at repeated roots. Here the balanced factors supply the missing
orders of vanishing. Indeed,
\[
 \Phi_k(t;Q_k)=(t^2-1)^{a-k}S_a(t),
\]
where
\[
 S_a(t)=(t+1)^{L-a}(t-3)^aH(t-2)
        +(t-1)^{L-a}(t+3)^aH(t+2).
\]
The equations at the distinct roots of $H$ give $S_a(z_j)=0$, hence
$H\mid S_a$. Consequently $Q_k\mid\Phi_k(t;Q_k)$ for every $k$.
The local algebras $J_k$ are nevertheless different from the original
ordered Bethe algebras; depermutation is essential in relating their
lengths.

\paragraph{Normalization at a repeated singular configuration.}
The $TQ$ interpretation of depermutation gives
\begin{equation}\label{eq:balanced-local-normalization}
 c_k(p)=\frac{(a!)^2}{n!}\nu_k(p),\qquad 0\le k\le a.
\end{equation}
To check the factor, there are $\binom ak^2$ compatible choices of
fixed labels. The remaining ordered root map has local degree
$((a-k)!)^2$, and the averaging factor in
\eqref{eq:reduced-fixed-pair-depermutation} gives
\[
 \frac{(k!)^2(n-2k)!}{n!}\binom ak^2
 \frac{((a-k)!)^2}{(n-2k)!}=\frac{(a!)^2}{n!}.
\]
Thus the required vanishing is precisely
\begin{equation}\label{eq:balanced-remainder-length-identity}
 \sum_{k=0}^a(-1)^k\binom Lk\nu_k(p)=0.
\end{equation}

\subsection{All roots in the singular clusters: $n=2a$}

\begin{proposition}\label{prop:pure-balanced-pairs}
Let $a\ge2$, $n=2a$, $L\ge4a$, and $p=(1^a,-1^a)$. Then
\begin{equation}\label{eq:pure-balanced-local-lengths}
 \nu_k(p)=
 \begin{cases}
 \displaystyle\binom{L-k}{a-k},&L\equiv a\pmod2,\\[4pt]
 \displaystyle\binom{L-k+1}{a-k},&L\not\equiv a\pmod2.
 \end{cases}
\end{equation}
In particular, $\operatorname{coeff}_p Z^{\mathrm{cand}}_{L,2a}=0$
for both parities of $L$.
\end{proposition}

\begin{proof}
For $k=a$ there are no free variables and $\nu_a=1$. Otherwise put
$s=a-k\ge1$ and $M=L-k$. Factor the variable polynomial locally as
\[
 Q(t)=A(t-1)B(t+1),\qquad
 A(u)=u^s-e_1u^{s-1}+\cdots+(-1)^se_s,
\]
with $B$ also monic of degree $s$. The two cluster factors at $p$ are
coprime, so their coefficients are local coordinates on coefficient
space. Divisibility by $Q$ is equivalent to the congruences
\[
 B(u)+u^MU(u)\equiv0\pmod{A(u)},\qquad
 A(v)+v^MV(v)\equiv0\pmod{B(v)},
\]
where
\[
 \begin{aligned}
 U(u)&=\frac{(u+4)^kA(u+2)B(u+4)}
                 {(u+2)^M(u-2)^kA(u-2)},\\
 V(v)&=\frac{(v-4)^kA(v-4)B(v-2)}
                 {(v-2)^M(v+2)^kB(v+2)}.
 \end{aligned}
\]
All denominators are units in the corresponding completed cluster
algebras. The first congruence solves uniquely for the coefficients
of $B$ in terms of those of $A$.

Give $u,v$ weight one and $e_j$ weight $j$. If $b_j$ denotes the
$j$th elementary coefficient of $B$, weighted division in the first
congruence gives
\[
 b_j-e_j=O_{\mathrm w}(M-s+j),\qquad 1\le j\le s.
\]
Since $M-s=L-a\ge3a$, these differences do not affect the constant
and linear weighted terms of the units. Substituting $B=A$ to these
orders gives, with $C=2^{2a-L}$ and $\alpha=(5a-2L)/4$,
\[
 \begin{aligned}
 U(u)&=(-1)^aC\left(1+\alpha u-\frac54e_1
                               +O_{\mathrm w}(2)\right),\\
 V(u)&=(-1)^LC\left(1-\alpha u+\frac54e_1
                               +O_{\mathrm w}(2)\right).
 \end{aligned}
\]
After eliminating $B$, the initial equations are therefore the
coefficients of $\operatorname{rem}_A u^M$ when $L\equiv a\pmod2$,
and, up to a common nonzero scalar, those of
\begin{equation}\label{eq:balanced-resonant-initial-polynomial}
 \operatorname{rem}_A
       \left[u^M\bigl((5a-2L)u-5e_1\bigr)\right]
\end{equation}
when the parities differ. Corrections from $B-A$ have strictly higher
weight, again because $M-s=L-a\ge3a$.

Both initial systems have only the common zero $e_1=\cdots=e_s=0$.
For the first, $A\mid u^M$ forces $A=u^s$. For the second, any
nonzero root of $A$ must be the unique simple root of
$(5a-2L)u-5e_1$. Thus, if such a root $\rho$ existed, we would have
$A=u^{s-1}(u-\rho)$ and $e_1=\rho$. Evaluation at $\rho$ would give
\[
 \bigl(5(a-1)-2L\bigr)\rho^{M+1}=0,
\]
which is impossible for $L\ge4a$. Hence $A=u^s$ here as well.

The initial equations consequently form a zero-dimensional weighted
complete intersection. Their weights are
$M-s+1,\ldots,M$ in the first case, and
$M-s+2,\ldots,M+1$ in the second. Higher-weight terms preserve the
local length. Dividing the product of these equation weights by
$1\cdot2\cdots s$ gives respectively $\binom Ms$ and
$\binom{M+1}s$, proving~\eqref{eq:pure-balanced-local-lengths}.

Finally,
\[
 \sum_{k=0}^a(-1)^k\binom Lk\binom{L-k}{a-k}
 =\binom La(1-1)^a=0.
\]
For the other parity, Pascal's identity adds
\[
 \sum_{k=0}^{a-1}(-1)^k\binom Lk\binom{L-k}{a-k-1}
 =\binom L{a-1}(1-1)^{a-1}=0,
\]
since $a\ge2$. Together with
\eqref{eq:balanced-local-normalization}, this proves the vanishing
of the corrected weight.
\end{proof}

\subsection{Additional roots: the nonresonant case}

The other roots cannot be $\pm3$, and no two of them differ by $2$:
the finite-chain argument used for one singular pair applies to the
distinct root values, regardless of the multiplicities at $\pm1$.
In particular, the following expression is well-defined:
\begin{equation}\label{eq:balanced-resonance-function}
 \Xi(p)=(-1)^a\frac{H(3)}{H(-1)}
           +(-1)^L\frac{H(-3)}{H(1)}.
\end{equation}
We call the configuration nonresonant when $\Xi(p)\ne0$.

\begin{proposition}\label{prop:balanced-nonresonant}
Assume $L\ge2n$, the relevant local algebras have finite length, and
$\Xi(p)\ne0$. Then
\[
 \nu_k(p)=\binom{L-k}{a-k}\nu_a(p),\qquad 0\le k\le a.
\]
Consequently the corrected weight at $p$ is zero for every $a\ge1$.
No simplicity assumption is needed for the remaining-root system.
\end{proposition}

\begin{proof}
For $k<a$, again put $s=a-k$ and $M=L-k$. Use the local factorization
\[
 Q(t)=A(t-1)B(t+1)H_y(t),\qquad
 H_y(t)=\prod_{j=1}^r(t-y_j),
\]
with $A,B$ monic of degree $s$ and $y$ near $z$. The cluster
congruences are those in the preceding proof, with their units
multiplied respectively by
$H_y(u+3)/H_y(u-1)$ and $H_y(v-3)/H_y(v+1)$.
After eliminating $B$, give the coefficients $e_j$ of $A$ weight $j$
and keep $y-z$ as parameters of weight zero. The constant weighted
term of $U+V$ is
\[
 2^{2a-L}\left[(-1)^a\frac{H_y(3)}{H_y(-1)}
                  +(-1)^L\frac{H_y(-3)}{H_y(1)}\right],
\]
which is a unit near $p$. Thus the cluster initial ideal is generated
by the coefficients of $\operatorname{rem}_A u^M$.

Setting all cluster coefficients to zero in the remaining equations
gives exactly the reduced $TQ$ system $J_a$ for $H_y$. Hence the
weighted initial local algebra is the tensor product of the local
algebra of $J_a$ at $H$ and
\[
 \C[e_1,\ldots,e_s]\big/
 \bigl(\text{coefficients of }\operatorname{rem}_A u^M\bigr).
\]
Both factors have finite length. The weighted complete-intersection
deformation therefore preserves length, and the second factor has
length $\binom Ms$ by the preceding calculation. This proves the
formula for $\nu_k$; for $k=a$ it is immediate. Finally,
\[
 \sum_{k=0}^a(-1)^k\binom Lk\nu_k(p)
 =\nu_a(p)\binom La(1-1)^a=0.
\]
Apply~\eqref{eq:balanced-local-normalization} to obtain the result.
\end{proof}

\subsection{One additional root: $n=2a+1$}

\begin{proposition}\label{prop:balanced-one-additional-root}
Let $a\ge2$, $n=2a+1$, and $L\ge4a+2$. The corrected weight is zero
at every configuration $p=(1^a,-1^a,z)$ with $z\ne\pm1$.
If $\Xi(p)=0$, its local remainder lengths are
\[
 \nu_k(p)=\binom{L-k+1}{a-k},\qquad 0\le k\le a.
\]
\end{proposition}

\begin{proof}
The nonresonant case is Proposition~\ref{prop:balanced-nonresonant}.
Suppose $\Xi(p)=0$. With $H(t)=t-z$, this condition reads
\[
 (-1)^a(z-3)(z-1)+(-1)^L(z+3)(z+1)=0.
\]
Thus $z=0$ when the parities differ, and $z^2=-3$ when they agree.
The remaining-root equation is
\[
 g_a(z)=(z+1)^{L-a}(z-3)^a-(z-1)^{L-a}(z+3)^a=0.
\]
Define
\[
 R(z)=-\frac2{z^2-1},\qquad S(z)=-\frac6{z^2-9},\qquad
 D(z)=(L-a)R(z)-aS(z).
\]
The logarithmic derivative of the ratio of the two terms in $g_a$
is $D$. At the possible resonant roots,
\[
 \begin{array}{c|c|c}
 &R+S&D\\ \hline
 z=0&8/3&\dfrac23(3L-4a)\\[4pt]
 z^2=-3&1&\dfrac{L-2a}{2}
 \end{array}
\]
so $g_a$ has a simple zero there and $\nu_a(p)=1$.

For $k<a$, put $s=a-k$ and $M=L-k$, and factor the variable
polynomial as $Q=A(t-1)B(t+1)(t-y)$. Eliminate $B$ as before.
The equation at $y$ also solves for $y$ as a function of the cluster
coefficients, with
\[
 y=z-\frac{R+S}{D}\,e_1+O_{\mathrm w}(2).
\]
This follows by expanding its logarithmic form: its linear terms are
$D(y-z)+(R+S)e_1$. Substituting into the cluster units, their
constant sum vanishes and their first weighted term, up to a nonzero
scalar, is $\alpha u+\beta e_1$, where
\[
 \alpha=\frac{5a-2L}{2}+R+S,\qquad
 \beta=-\frac52+\frac{(R+S)^2}{D}.
\]
The displayed values and $L\ge4a+2$ give $\alpha<0$ and $\beta<0$.
In particular, $\alpha\ne0$ and $\alpha+\beta\ne0$.

The cluster initial equations are therefore the coefficients of
\[
 \operatorname{rem}_A\bigl[u^M(\alpha u+\beta e_1)\bigr].
\]
They have only the zero solution in coefficient space. Indeed, any
nonzero root $\rho$ of $A$ would be the unique simple nonzero root
of $\alpha u+\beta e_1$, forcing $e_1=\rho$ and
$(\alpha+\beta)\rho=0$, a contradiction. The weighted
complete-intersection calculation from
Proposition~\ref{prop:pure-balanced-pairs} now gives
$\nu_k=\binom{M+1}s$. The same two binomial identities prove
\eqref{eq:balanced-remainder-length-identity}, hence the required
zero weight.
\end{proof}

\subsection{A jet-space Koszul model without deforming the equations}
\label{sec:jet-koszul-model}

We now assume that the remaining roots are distinct and that the fully
reduced zero is simple. Write
\[
 A_k=\mathcal O_{V(J_k),Q_k},\qquad \nu_k=\dim_\C A_k,
 \qquad A_a=\C.
\]
The construction below uses the completed local equations at the
homogeneous point. It does not use a deformation of the Bethe equations
or a spectral comparison. This particular Koszul model depends on an
explicit weighted initial-form hypothesis, proved below on a locus
described by the reduced Jacobian. Section~\ref{sec:direct-residual-ideal}
then gives a residual-length interpretation without that weighted
hypothesis, using an auxiliary intersection argument.

\paragraph{The residual one-pair multiplicity.}
Put $m=L-a+1$. In the system $J_{a-1}$ one singular pair remains free.
The cluster factorization and simplicity of $A_a$ allow the ordinary
roots and one member of this pair to be eliminated. As in the one-pair
calculation, the remaining equation is $u^m h(u)=0$: the factor $u^m$
comes from the vacuum exponent $L-(a-1)$. Thus
\begin{equation}\label{eq:jet-residual-order}
 A_{a-1}\simeq\C[[u]]/(u^m h(u)),\qquad
 q=\ord_0h\ge0,\qquad D=m+q=\nu_{a-1}.
\end{equation}
Finiteness ensures $h\ne0$; a unit has order zero. In particular,
$A_{a-1}\simeq\C[[u]]/(u^D)$. The integer $q$ is defined by the actual
one-pair local equations, before any hypothesis about the other lengths.

For $k<a$, set $s=a-k$ and factor the variable polynomial near $Q_k$ as
\[
 Q(t)=C_s(t-1)B_s(t+1)H_y(t),\qquad
 C_s(u)=u^s-c_1u^{s-1}+\cdots+(-1)^sc_s.
\]
The first cluster congruence solves for $B_s$, and simplicity of $A_a$
solves the ordinary-root equations for $y$. Hence
\[
 A_k\simeq\C[[c_1,\ldots,c_s]]/(G_{k,1},\ldots,G_{k,s}).
\]
Give $c_j$ weight $j$ and take initial forms of smallest weight.

\begin{definition}[Weighted local hypothesis]\label{def:jet-weighted-hypothesis}
With $D$ determined by~\eqref{eq:jet-residual-order}, suppose that, for
every $k<a$, the generators can be chosen so that their weighted
initial forms are a regular sequence of degrees
\begin{equation}\label{eq:jet-initial-degrees}
 D,D+1,\ldots,D+s-1,\qquad s=a-k.
\end{equation}
For $s=1$ this follows from~\eqref{eq:jet-residual-order}; the additional
condition concerns the higher collision strata.
\end{definition}

This is a condition on leading local equations, stronger than simplicity
of $A_a$. We do not assert that it holds at every Bethe configuration.
In particular, it must not be replaced by an unproved assertion that
$q$ is always zero or one.

\begin{proposition}[Conditional jet-space index]\label{prop:jet-koszul-index}
Assume the weighted local hypothesis. Then
\begin{equation}\label{eq:jet-local-lengths}
 \nu_k=\binom{L+q-k}{a-k},\qquad 0\le k\le a.
\end{equation}
There is an explicit finite complex of vector spaces with term
dimensions $\binom Lk\nu_k$, no positive homology, and degree-zero
homology $M_p$ satisfying
\begin{equation}\label{eq:jet-koszul-index}
 \boxed{\sum_{k=0}^a(-1)^k\binom Lk\nu_k
       =\dim_\C M_p=\binom qa.}
\end{equation}
Here $\binom qa=0$ for $q<a$. In particular, $q<a$ implies the desired
balanced cancellation; $q\le1$ suffices when $a\ge2$.
\end{proposition}

\begin{proof}
The filtered Koszul argument for a regular sequence of initial forms
identifies the associated graded quotient with that initial complete
intersection. Its Hilbert series and length are
\[
 \frac{\prod_{j=1}^s(1-z^{D+j-1})}
      {\prod_{j=1}^s(1-z^j)},\qquad
 \nu_k=\prod_{j=1}^s\frac{D+j-1}{j}
       =\binom{D+s-1}s.
\]
This proves~\eqref{eq:jet-local-lengths}, including $\nu_a=1$.

Let $W=\C[u]/(u^D)$ as a vector space and $S=\Sym_\C(W)$, with all
elements of $W$ placed in polynomial degree one. Define the linear map
\[
 \lambda:\C^L\longrightarrow W,\qquad
 \lambda(\varepsilon_i)=[u^{i-1}],\quad 1\le i\le L.
\]
In this formula $[1]\in W$ is a degree-one generator of $S$, not the
scalar unit of $S$. The map has rank $\min(L,D)$. Take the internal
degree-$a$ part of the ordinary Koszul complex on these $L$ linear
forms:
\[
 C_k=S_{a-k}\otimes\bigwedge^k\C^L,\qquad
 d(f\otimes\varepsilon_{i_1}\wedge\cdots\wedge\varepsilon_{i_k})
 =\sum_{j=1}^k(-1)^{j-1}[u^{i_j-1}]f\otimes
       \varepsilon_{i_1}\wedge\cdots\widehat{\varepsilon_{i_j}}
                      \cdots\wedge\varepsilon_{i_k}.
\]
Its term dimensions are
\[
 \dim C_k=\binom Lk\dim\Sym^{a-k}W
          =\binom Lk\binom{D+a-k-1}{a-k}
          =\binom Lk\nu_k.
\]

If $D\ge L$, the linear forms are independent and form a regular
sequence. Thus positive homology vanishes and
\[
 M_p=H_0(C_\bullet)=\Sym^a(\operatorname{coker}\lambda).
\]
For $D>L$, the cokernel has dimension $D-L=q-a+1$, so this space has
dimension $\binom qa$; for $D=L$ it is zero, since $a\ge1$.

If $D<L$, split the domain into a complement mapping isomorphically
onto $W$ and $\ker\lambda$. The full Koszul homology is then
$\bigwedge^j\ker\lambda$ in internal degree $j$. However,
\[
 \dim\ker\lambda=L-D=a-q-1<a.
\]
Its internal degree-$a$ part is therefore zero in every homological
degree. This also gives $M_p=0$ and $\binom qa=0$. Taking the Euler
characteristic proves~\eqref{eq:jet-koszul-index}.
\end{proof}

This realizes the index as a dimension under the weighted hypothesis.
It does not construct canonical maps between the actual algebras $A_k$:
the terms are symmetric powers of a jet space with the required
dimensions. In particular, the differentials are not $A_0$-linear maps
between $A_k^{\binom Lk}$, so the socle obstruction discussed in
Section~\ref{sec:normalized-local-complex} does not apply. For $a=1$
the index is $q$, agreeing with the residual one-pair length; for $a=2$
it is $q(q-1)/2$.

\paragraph{An explicit first-weight criterion.}
The weighted hypothesis can be proved without spectral input under
nonvanishing conditions expressed in the fully reduced Jacobian.
For the ordinary roots set
\[
 \mathcal B_j(y)=
 \frac{(y_j+1)^{L-a}(y_j-3)^a}
      {(y_j-1)^{L-a}(y_j+3)^a}
 \prod_{\ell\ne j}\frac{y_j-y_\ell-2}{y_j-y_\ell+2},\qquad
 \mathcal J_{j\ell}=\partial_{y_\ell}\log\mathcal B_j(z).
\]
These logarithmic derivatives are rational functions; no choice of a
logarithm is needed. At the solution $\mathcal B_j(z)=1$, and simplicity
of the fully reduced system makes $\mathcal J$ invertible. Define
\begin{equation}\label{eq:jet-alpha-beta}
 \begin{aligned}
 \gamma_j&=\frac1{z_j+1}+\frac1{z_j+3}
                 -\frac1{z_j-1}-\frac1{z_j-3},\\
 \alpha&=\frac{5a-2L}{2}+\sum_j\gamma_j,\qquad
 \beta=-\frac52+\gamma^{\mathsf T}\mathcal J^{-1}\gamma.
 \end{aligned}
\end{equation}
The transpose here is algebraic, not a conjugate transpose. For no
ordinary roots the sum and quadratic expression are zero.

\begin{proposition}\label{prop:jet-first-weight-criterion}
Retain the finiteness and simplicity assumptions above.
If $\Xi(p)\ne0$, the weighted hypothesis holds with $q=0$.
If
\begin{equation}\label{eq:jet-resonant-open-conditions}
 \Xi(p)=0,\qquad \alpha\ne0,\qquad \alpha+\beta\ne0,
\end{equation}
it holds with $q=1$. Thus the conditional index theorem applies on
these loci, and its index vanishes for every $a\ge2$ there.
\end{proposition}

\begin{proof}
Put $s=a-k$ and $M=L-k$. The cluster congruences take the form
\[
 B_s(u)+u^MU(u)\equiv0\pmod{C_s(u)},\qquad
 C_s(v)+v^MV(v)\equiv0\pmod{B_s(v)}.
\]
The units are those in the proof of
Proposition~\ref{prop:pure-balanced-pairs}, multiplied respectively by
$H_y(u+3)/H_y(u-1)$ and $H_y(v-3)/H_y(v+1)$.
Eliminating $B_s$ changes its $j$th elementary coefficient from $c_j$
only in weights at least $M-s+j=L-a+j$. These changes do not affect
the constant or first weighted terms. The equations at the ordinary
roots have linear terms
\[
 \mathcal J(y-z)+\gamma c_1=0,
 \qquad y=z-\mathcal J^{-1}\gamma c_1+O_{\mathrm w}(2).
\]
Both statements are independent of $k$, since the fixed and variable
cluster factors have total multiplicity $a$ at each sign.

The constant term of $U+V$ is $2^{2a-L}\Xi(p)$. If it is nonzero,
the initial equations are the coefficients of
$\operatorname{rem}_{C_s}u^M$. Their degrees are
$M-s+1,\ldots,M$, and they have only the zero solution. This proves
the $q=0$ assertion, also recovering the simple-reduced case of
Proposition~\ref{prop:balanced-nonresonant}.

Suppose $\Xi(p)=0$. The nonzero constant terms of $U,V$ are opposite.
After division by the first of these constants, their first weighted
sum is
\[
 \alpha u+\beta c_1.
\]
For completeness, the pure-cluster contributions to the difference of
their logarithmic derivatives are $(5a-2L)/2$ in $u$ and $-5/2$ in
$c_1$. The $H_y$ factors add $\sum_j\gamma_j$ to the former. Their
relative derivative in $y$ is $-\gamma^{\mathsf T}$, so substitution
of $y-z=-\mathcal J^{-1}\gamma c_1+O_{\mathrm w}(2)$ adds
$\gamma^{\mathsf T}\mathcal J^{-1}\gamma$ to the latter. This gives
exactly~\eqref{eq:jet-alpha-beta}.

Consequently the initial equations are the coefficients of
\[
 \operatorname{rem}_{C_s}
       \bigl[u^M(\alpha u+\beta c_1)\bigr].
\]
They have only the zero solution under
\eqref{eq:jet-resonant-open-conditions}. Indeed, any nonzero root
$\rho$ of $C_s$ must be the unique simple nonzero root of the linear
factor, forcing $C_s=u^{s-1}(u-\rho)$ and $c_1=\rho$. This would give
$(\alpha+\beta)\rho=0$, a contradiction. The initial sequence is
therefore regular, with degrees $M-s+2,\ldots,M+1$. When $s=1$ its
leading equation is $(\alpha+\beta)c_1^{M+1}$, so the residual order
in~\eqref{eq:jet-residual-order} is $q=1$. The same degrees for all
$s$ are precisely~\eqref{eq:jet-initial-degrees} with
$D=L-a+2$, proving the assertion.
\end{proof}

\paragraph{The remaining exceptional loci.}
The coefficient $\alpha$ is the negative of the second residue
obstruction $K$ in~\eqref{eq:balanced-second-obstruction}.
Thus the criterion also computes every individual length on the
simple-reduced part of the locus $\Xi=0$, $K\ne0$ where
$\alpha+\beta\ne0$. It recovers the two pure-cluster parities and
the resonant one-additional-root calculation above.
The unresolved first-weight cases are
\[
 \Xi=0,\quad\alpha=0,
 \qquad\text{or}\qquad
 \Xi=0,\quad\alpha+\beta=0.
\]
In the first case with $\beta\ne0$ and $s\ge2$, the candidate initial
equations vanish on $c_1=0$ and are not isolated, even though the
one-pair residual order is still $q=1$. In the second case a nonzero
root $\rho$ is no longer excluded by the initial divisibility equation,
and the one-pair residual order is at least two. Higher weighted terms
must be analyzed at these loci; simplicity of $A_a$ does not justify
discarding them. Neither the weighted hypothesis nor a bound $q<a$
has been proved there.

\subsection{A residual ideal from the homogeneous cluster equations}
\label{sec:direct-residual-ideal}

Under the same simple-reduced assumption, the alternating sum has a
meaning without the weighted hypothesis: it is the length of an
explicit residual ideal. The following general local intersection
identity is the main ingredient. Its proof separates a fixed divisor
in an auxiliary space; no twist, site inhomogeneities, or transfer
matrix is introduced.

\begin{proposition}[Local divisor subtraction]
\label{prop:local-divisor-subtraction}
Let $L\ge a\ge1$ and
\[
 C(t)=t^a-c_1t^{a-1}+\cdots+(-1)^ac_a.
\]
Let $W(t;c)$ be a convergent power series near $(t,c)=(0,0)$.
Suppose the ideal of coefficients of
$\operatorname{rem}_C(t^LW(t;c))$ has finite colength at $c=0$.
For $0\le k\le a$, write $C=t^kC_s$, $s=a-k$, and let $\mu_k$
be the local length of
\[
 \operatorname{rem}_{C_s}
 \bigl(t^{L-k}W(t;t^kC_s)\bigr)=0
 \quad\text{at }C_s=t^s.
\]
Here $W(t;t^kC_s)$ denotes substitution of the coefficients of $t^kC_s$
for $c$, and $\mu_a=1$. Then the residual ideal
\[
 I_{\mathrm{res}}=
 \bigl(\text{coefficients of }\operatorname{rem}_C W(t;c)\bigr)
 \subseteq\C\{c_1,\ldots,c_a\}
\]
has finite colength, allowing the unit ideal, and
\begin{equation}\label{eq:local-divisor-subtraction}
 \boxed{\sum_{k=0}^a(-1)^k\binom Lk\mu_k
        =\length\bigl(\C\{c\}/I_{\mathrm{res}}\bigr).}
\end{equation}
The same lengths are obtained in the completed local rings.
\end{proposition}

\begin{proof}
The residual ideal contains the original remainder ideal: multiplication
by $t^L$ in the free rank-$a$ algebra $\C\{c\}[t]/(C)$ expresses the
original remainder coefficients as linear combinations of the residual
ones. Thus the residual quotient is finite. Likewise, restriction to
$C=t^kC_s$ identifies the original remainder with $t^k$ times the
smaller remainder in the statement. Hence all the $\mu_k$ are finite.

Here is a local intersection proof of the identity. Replace the fixed
factor $t^L$ by
\[
 P_\delta(t)=\prod_{i=1}^L(t-\delta_i)
\]
and replace $W$ by $W+h_b$, where $h_b(t)=\sum_{j=0}^{a-1}b_jt^j$.
For $I\subseteq[L]$, $|I|=k\le a$, put
$P_I=\prod_{i\in I}(t-\delta_i)$ and impose $C=P_IC_s$. The smaller
remainder equations are
\[
 \operatorname{rem}_{C_s}
 \left(P_{[L]\setminus I}(t)
             \bigl(W(t;P_IC_s)+h_b(t)\bigr)\right)=0.
\]
Their central fibers have lengths $\mu_k$. They have $s$ equations in
$s$ coefficient variables, so conservation of isolated complete-
intersection multiplicity counts $\mu_k$ nearby points for generic
small $(\delta,b)$. The same applies to the residual equations
$C\mid W+h_b$.

We justify choosing the parameters so that these counts are ordinary
counts of simple points. Take distinct $\delta_i$ and stratify by the
exact subset $T$ of their roots contained in $C$. Write $C=P_TH$,
initially allowing $H$ to contain selected nodes again, but requiring
it to avoid the nodes outside $T$. Since $P_{[L]\setminus T}$ is a
unit modulo $H$, the remaining condition is
\[
 H\mid W(t;P_TH)+h_b(t).
\]
Variation of $h_b$ maps surjectively to $\C[t]/(H)$, including when
$H$ has repeated roots. The residual incidence scheme is therefore
smooth and projects smoothly onto the coefficients of $H$ and the
$\delta_i$. Repeated roots of $H$, or a root of $H$ at a selected node,
form proper hypersurfaces in this coefficient space and are absent
for generic parameters. Thus repeated site roots are excluded as
well, rather than being omitted from the stratification.
For a selected site $i\in T$, the extra condition
$(W+h_b)(\delta_i)=0$ is also avoidable: evaluation there and the
remainder modulo $H$ are independent conditions, with
$\deg H+1\le a$. The finite projections to the parameter space are
generically etale. Thus, simultaneously for the finitely many strata,
the residual roots are distinct, and the factors selected from
$P_\delta$ are transverse simple roots of the product. The same
argument covers $T=\varnothing$, whose residual polynomial is
independent of $\delta$, so its roots avoid generic $\delta_i$.

At any such nearby full solution, let $T$ be its exact selected subset.
It is counted in the $I$-system precisely when $I\subseteq T$.
Its total alternating contribution is
\[
 \sum_{I\subseteq T}(-1)^{|I|}
 =\begin{cases}1,&T=\varnothing,\\0,&T\ne\varnothing.\end{cases}
\]
The surviving points are exactly the nearby residual solutions
$C\mid W+h_b$. Conservation of their local intersection numbers gives
\eqref{eq:local-divisor-subtraction}. This argument is an auxiliary
separation of the fixed divisor, not an assertion about deformed Bethe
solutions.
\end{proof}

\begin{proposition}[Compatible scalar cluster equation]
\label{prop:scalar-cluster-equation}
Assume $L\ge2n$, the standing local finiteness hypotheses, distinct
ordinary roots, and simplicity of the fully reduced zero $A_a$.
There is a convergent germ $W(u;c_1,\ldots,c_a)$ such that, after
eliminating the negative cluster and ordinary roots,
\begin{equation}\label{eq:scalar-full-cluster}
 A_0\simeq\C[[c_1,\ldots,c_a]]/
       \bigl(\text{coefficients of }\operatorname{rem}_C u^LW\bigr).
\end{equation}
For every $k$, the same germ restricted to $C=u^kC_s$ presents $A_k$
by the coefficients of
\begin{equation}\label{eq:scalar-reduced-cluster}
 \operatorname{rem}_{C_s}
          \bigl(u^{L-k}W(u;u^kC_s)\bigr),\qquad s=a-k.
\end{equation}
\end{proposition}

\begin{proof}
Factor the full polynomial as $Q(t)=C(t-1)B(t+1)H_y(t)$, with $C,B$
monic of degree $a$. Its cluster congruences are
\[
 B(u)+u^LU(u)\equiv0\pmod{C(u)},\qquad
 C(u)+u^LV(u)\equiv0\pmod{B(u)},
\]
where the following are analytic units near the origin:
\[
 \begin{aligned}
 U(u)&=\frac{C(u+2)B(u+4)H_y(u+3)}
              {(u+2)^LC(u-2)H_y(u-1)},\\
 V(u)&=\frac{C(u-4)B(u-2)H_y(u-3)}
              {(u-2)^LB(u+2)H_y(u+1)}.
 \end{aligned}
\]
The first remainder equations solve uniquely for $B$; their Jacobian
in its coefficients is the identity at the origin. The ordinary-root
equations then solve uniquely for $y$ by simplicity of the fully
reduced system. Henceforth $B,y,U,V$ are analytic functions of $c$.

The first congruence gives an exact Weierstrass division identity
\[
 B+u^LU=CE.
\]
Divide $C+u^LV$ by $B$ and write
\[
 C+u^LV=BF+R_B,\qquad \deg_uR_B<a.
\]
Both quotient germs $E,F$ have constant term one. Substituting the
first identity into the second gives
\[
 R_B=C(1-EF)+u^L(V+UF).
\]
Taking the remainder modulo $C$ leaves $R_B$ itself, since its degree
is less than $a$. Therefore, with
\begin{equation}\label{eq:effective-cluster-W}
 W=V+UF,
\end{equation}
we obtain the exact identity $R_B=\operatorname{rem}_C(u^LW)$.
This proves~\eqref{eq:scalar-full-cluster} without discarding any
higher weighted terms.

It remains to check compatibility, which is essential for subtraction.
Restrict to $C=u^kC_s$. The identity $B+u^LU=CE$ and $L\ge k$ force
$B=u^kB_s$. The restrictions of the units $U,V$ are exactly the units
of the reduced $k$-pair system: substituting $C=u^kC_s$, $B=u^kB_s$
in their displayed formulas produces its factors $(u+4)^k,(u-2)^k$
and their negative-cluster analogues. The ordinary equations also
agree after cancelling units at the ordinary roots. Uniqueness in
the implicit-function theorem identifies all eliminated functions
with their reduced-system counterparts.

Finally, dividing the two Weierstrass identities by $u^k$ gives
\[
 B_s+u^{L-k}U=C_sE,\qquad
 C_s+u^{L-k}V=B_sF+R_s.
\]
Uniqueness of monic division gives $R_B=u^kR_s$ and leaves the same
quotients $E,F$, hence the same restricted $W$. Taking the remainder
modulo $C_s$ proves~\eqref{eq:scalar-reduced-cluster}.
\end{proof}

\begin{proposition}[Residual meaning of the alternating index]
\label{prop:direct-residual-index}
Under the hypotheses of Proposition~\ref{prop:scalar-cluster-equation},
define the finite residual algebra
\begin{equation}\label{eq:direct-residual-algebra}
 \mathcal R_p=\C[[c_1,\ldots,c_a]]/
       \bigl(\text{coefficients of }\operatorname{rem}_C W(u;c)\bigr).
\end{equation}
The zero algebra is allowed. Then
\begin{equation}\label{eq:direct-residual-index}
 \boxed{\sum_{k=0}^a(-1)^k\binom Lk\nu_k(p)
                  =\length\mathcal R_p\ge0.}
\end{equation}
Consequently its corrected ordered weight is
\[
 \operatorname{coeff}_p Z^{\mathrm{cand}}_{L,n}
       =\frac{(a!)^2}{n!}\length\mathcal R_p.
\]
No weighted initial-form or spectral assumption is required.
\end{proposition}

\begin{proof}
Apply Proposition~\ref{prop:local-divisor-subtraction} to the compatible
scalar equations in Proposition~\ref{prop:scalar-cluster-equation}.
The local lengths $\mu_k$ in that proposition are exactly $\nu_k(p)$.
The weight formula follows from~\eqref{eq:balanced-local-normalization}.
\end{proof}

The presentation of $\mathcal R_p$ uses the indicated choices of local
coordinates and elimination; its length is independent of those choices
by~\eqref{eq:direct-residual-index}. For one pair it recovers the
residual algebra $\C[[u]]/(h)$. Under the weighted hypothesis its length
is $\binom qa$, so the jet-space Koszul model computes this residual
length. Without that hypothesis, $q$ alone need not determine the index.
For example, in the abstract divisor problem with $a=2$ and
$W(u;c)=u^2+c_1$, the residual ideal is $(c_1-c_2,c_1)$ and has length
one, although the one-pair residual order is $q=1$. This is not asserted
to be a Bethe configuration; it shows why the weighted hypothesis cannot
be inferred solely from that order.

At a Bethe point the constant term of $W$ is $2^{2a-L}\Xi(p)$.
When $\Xi=0$, its first weighted part is a nonzero scalar times
$\alpha u+\beta c_1$, with $\alpha,\beta$ as above, because $F-1$ has
weight at least $L-a$. Thus $\Xi\ne0$ makes the residual ideal the unit
ideal. If $a\ge2$, $\Xi=0$ and $\alpha\ne0$ also make it the unit
ideal: the coefficient of $u$ in $\operatorname{rem}_C W$ has nonzero
constant term. This proves cancellation there even without assuming
$\alpha+\beta\ne0$; that extra condition was needed for the explicit
length pattern, not for residual vanishing.

The remaining balanced problem has now become a concrete local ideal
problem: show that $\mathcal R_p=0$ for $a\ge2$. The exceptional case
$\Xi=\alpha=0$ is not settled by the first weighted terms. The argument
does not prove that a nonzero residual algebra has length one, nor remove
the simple-reduced hypothesis. Its only deformation is the auxiliary
divisor separation in the general intersection lemma; a completely
deformation-free homological proof of that lemma is not supplied here.

\subsection{A local geometric complex from normalized deformations}
\label{sec:normalized-local-complex}

There is a geometric realization of the alternating local length which
does not use the Bethe algebra or its spectrum. Its terms come from
normalizations of one-parameter deformations of the local schemes.
They have the required dimensions, but are not direct sums of the
original local algebras. This distinction is essential.

In this subsection $I\subseteq[L]$ labels sites of the chain, rather
than fixed root labels. For site parameters
$\theta=(\theta_1,\ldots,\theta_L)$ and
$I\subseteq[L]$, put $r=n-2|I|\ge0$ and define
\[
 \begin{aligned}
 D_i(t)&=(t-\theta_i)^2-1,\\
 a_I(t)&=\prod_{i\notin I}(t-\theta_i+1)
          \prod_{i\in I}(t-\theta_i-3),\\
 d_I(t)&=\prod_{i\notin I}(t-\theta_i-1)
          \prod_{i\in I}(t-\theta_i+3),\\
 \Phi_I(t;Q)&=a_I(t)Q(t-2)+d_I(t)Q(t+2).
 \end{aligned}
\]
Let $\mathcal X_I\to\A^L_\theta$ be the scheme defined by the
remainder of $\Phi_I$ modulo a monic degree-$r$ polynomial $Q$.
At $\theta=0$ this is the reduced remainder scheme $V(J_{|I|})$.
For $j\notin I$ with $2(|I|+1)\le n$, the identity
\begin{equation}\label{eq:inhomogeneous-factor-insertion}
 \Phi_I(t;D_jQ)=D_j(t)\Phi_{I\cup\{j\}}(t;Q)
\end{equation}
gives a closed immersion
$\mathcal X_{I\cup\{j\}}\hookrightarrow\mathcal X_I$.
Indeed, the remainder on the left is $D_j$ times the smaller
remainder, and multiplication by the monic polynomial $D_j$ does
not change the ideal generated by its coefficients.

\begin{proposition}[Finite flat remainder families]
\label{prop:geometric-remainder-flatness}
Assume $L\ge2n$. Each $\mathcal X_I\to\A^L_\theta$ is finite
flat of degree
\[
 N_{L,r}=\binom{L+r-2}{r}.
\]
There is a nonempty Zariski open set of site parameters on which all
these fibers are reduced. On this open set their polynomials $Q$
have distinct roots. Every pair of roots differing by $2$ is exactly
$\theta_j-1,\theta_j+1$ for some $j\notin I$, with both roots simple.
\end{proposition}

\begin{proof}
The case $r=0$ is immediate. For $r>0$, write
$Q=t^r+q_1t^{r-1}+\cdots+q_r$ and give $q_j$ weight $j$.
Set $b=L-4|I|\ge2r$. Taylor expansion shows that the weighted
initial equations of the remainder coefficients are the coefficients of
\[
 \operatorname{rem}_Q\left(4t^{L-2}
               \bigl(t^2Q''-btQ'\bigr)\right).
\]
Their respective weights are $L-1,L,\ldots,L+r-2$; they are
independent of $\theta$. To find their common zero set, factor
$Q=t^gH$ with $H(0)\ne0$. The divisibility condition implies
$H\mid t^2Q''-btQ'$. Since $t^g$ also divides this expression,
\[
 t^2Q''-btQ'=r(r-1-b)Q.
\]
The eigenvalues $m(m-1-b)$ for $0\le m\le r$ are distinct,
because $b\ge2r$. Hence $Q=t^r$.

The initial equations therefore form a regular sequence in the
weighted polynomial ring. The filtered Koszul complex gives the
associated graded quotient by these initial equations. It has length
\[
 \frac{(L-1)L\cdots(L+r-2)}{1\cdot2\cdots r}
 =N_{L,r}.
\]
The same initial equations work over $\C[\theta]$, giving a
uniform finite spanning set and this length in every fiber. The
finite complete-intersection family is consequently flat.

We give the generic-reducedness argument, including the boundary
where a site pair occurs. Choose distinct site parameters with no
differences in $2\mathbb Z\cap[-2n-6,2n+6]$. At the bottom and
top of a maximal string of roots spaced by $2$, evaluation of
$Q\mid\Phi_I$ forces a zero of
$a_I$ and of $d_I$, respectively. The nonresonance of the sites
then leaves only the pair $\theta_j-1,\theta_j+1$ at a positive
site $j\notin I$. The two multiplicities cannot both exceed one:
if they did, the vanishing order at the lower endpoint would be exactly
one, less than its multiplicity in $Q$.

Remove one copy of every such pair, and let $J\supseteq I$ be the
resulting set of negative sites. The remaining polynomial $H$ has
degree $s=n-2|J|$ and is coprime to its shifts by $2$ and to the
two vacuum polynomials of the $J$-system. Temporarily use the
coefficients of
\[
 W(t)=\prod_{j\notin J}(t-\theta_j),\qquad M=L-|J|,
\]
as the positive-site parameters, and keep the negative centers
$c_i=\theta_i$, $i\in J$. For fixed $H,c$, the remainder equations
are linear in the coefficients of the monic polynomial $W$:
\begin{equation}\label{eq:geometric-interpolation-equations}
 \prod_{i\in J}(t-c_i-3)W(t+1)H(t-2)
 +\prod_{i\in J}(t-c_i+3)W(t-1)H(t+2)
 \equiv0\pmod H.
\end{equation}
Since $H(t-1)$ and $H(t+1)$ are coprime, interpolation at their
roots, with multiplicities, prescribes the two classes
$\delta W(t-1),\delta W(t+1)$ independently. Here $M\ge2s$.
The linear map in~\eqref{eq:geometric-interpolation-equations}
is therefore surjective. The admissible incidence scheme is smooth
of dimension $L$, and its projection to $(H,c)$ is smooth. Its
projection to the site parameters is quasi-finite by the first part
of the proof, hence is generically etale in characteristic zero.
Furthermore, a repeated root of $H$, or a root $c_i\pm1$, cuts
out a proper hypersurface in the $(H,c)$ space. These conditions
are absent for generic site parameters. Thus the residual roots and
all the site-pair roots are generically distinct.

It remains to check transverse reducedness at an inserted pair.
Write its center as $c$, its roots as $c+1+u,c-1+v$, and the
product of all other root factors as $H_*$. In the system before
insertion write $a=(t-c+1)a_*$ and $d=(t-c-1)d_*$. The two
evaluation equations have linear parts
\[
 A(u-v)+Bu,\qquad Cv+D(v-u),
\]
where
\[
 \begin{aligned}
 A&=-4a_*(c+1)H_*(c-1),&
 B&=8d_*(c+1)H_*(c+3),\\
 C&=8a_*(c-1)H_*(c-3),&
 D&=-4d_*(c-1)H_*(c+1).
 \end{aligned}
\]
Their determinant is $AC+BC+BD$. It is not identically zero on
the residual incidence scheme. Indeed, in addition to the $2s$
interpolation conditions above, one can prescribe the $3|J|$
values $W(c_i),W(c_i-2),W(c_i+2)$ independently: their nodes are
distinct and avoid the roots of $H(t-1)H(t+1)$, and
\[
 M=L-|J|\ge2s+3|J|.
\]
Holding the other values fixed, the term $AC$ has a nonzero
coefficient of $W(c+2)W(c)$, whereas $BC$ and $BD$ do not
involve $W(c+2)$. Thus the determinant is generically nonzero.
The endpoint equations have zero derivatives in the residual root
directions and in other pairs' transverse directions. Their blocks,
together with the residual Jacobian, give an invertible full
Jacobian for generic parameters. There are finitely many $I,J$,
so a common open set works for all of them. Passing back from $W$
to its distinct ordered roots is etale and preserves these conclusions.
\end{proof}

\begin{proposition}[Normalized local complex]
\label{prop:normalized-local-complex}
Let $L\ge2n$ and let
$Q_0=(t^2-1)^aH$, with $H(\pm1)\ne0$, satisfy the full $TQ$
relation. Repetitions among the roots of $H$ are allowed. For
$Q_k=Q_0/(t^2-1)^k$, let $\nu_k(p)$ be the local remainder length.
A generic one-parameter deformation of the site parameters determines
a finite complex of vector spaces $C_\bullet(p)$ such that
\begin{equation}\label{eq:normalized-local-complex-dimensions}
 \dim C_k(p)=\binom Lk\nu_k(p)\quad(0\le k\le a),\qquad
 H_k(C_\bullet(p))=0\quad(k>0).
\end{equation}
Its degree-zero homology $M_p$ is the special fiber of the
normalization of the nearby components containing no exact site pair.
Consequently, without spectral input,
\begin{equation}\label{eq:geometric-local-positivity}
 \boxed{\sum_{k=0}^a(-1)^k\binom Lk\nu_k(p)
              =\dim_\C M_p\ge0.}
\end{equation}
The construction may depend on the chosen deformation; the displayed
dimension does not.
\end{proposition}

\begin{proof}
Choose a line $\theta_i=\varepsilon c_i$ whose generic point lies
in the common open set of Proposition~\ref{prop:geometric-remainder-flatness},
and work over $T=\C[[\varepsilon]]$. In $\mathcal X_\varnothing$
take the finite flat local factor specializing to $Q_0$. For each
$I\subseteq[L]$ with $2|I|\le n$, take its inverse image in $\mathcal X_I$ under
insertion of the factors $D_i$, $i\in I$, and denote its coordinate
algebra by $\mathcal A_I$. Henselian decomposition of a finite
algebra makes these local factors well-defined. For $|I|=k\le a$,
the special fiber is the local algebra of $J_k$ at $Q_k$, so
$\mathcal A_I$ is finite free of rank $\nu_k(p)$ over $T$.
For $k>a$ the inverse image is empty.

The generic fibers are reduced. Flatness implies that the total
algebras $\mathcal A_I$ are reduced as well. Let
$\mathcal B_I$ be their normalizations. These are finite products
of discrete valuation rings, finite free over $T$, with the same
ranks as $\mathcal A_I$; for the standard normalization facts see
\cite[Tag~035E]{Stacks}. Normalization separates the generic
components that can meet at the homogeneous point.

For $J\subset I$, factor insertion is a closed immersion whose
generic fiber selects a subset of the reduced points. Consequently
$\mathcal B_I$ is canonically a direct factor of $\mathcal B_J$.
There is an extension-by-zero map
\[
 \iota_{I,J}:\mathcal B_I\longrightarrow\mathcal B_J.
\]
These maps compose according to inclusions of subsets. Define
\[
 \mathcal C_k=\bigoplus_{|I|=k}\mathcal B_I,
 \qquad
 \left.d_k\right|_{\mathcal B_I}
 =\sum_{j\in I}(-1)^{\#\{i\in I:i<j\}}
             \iota_{I,I\setminus\{j\}}.
\]
The usual cancellation of signs gives $d_{k-1}d_k=0$.

To compute homology, decompose by the normal components of the
full local family. For one such component, let $S$ be the set of
sites whose factors $D_i$ divide its generic polynomial. Its
contribution to $\mathcal C_k$ is one copy of its normal ring for
each $k$-element subset of $S$. The differential is contraction
on the exterior algebra of the free module with basis $S$, sending
every basis vector to $1$. If $S\ne\varnothing$, exterior
multiplication by any chosen basis vector is a contracting homotopy.
If $S=\varnothing$, only degree zero remains. The whole complex
is therefore split exact in positive degrees, with degree-zero
homology equal to the product of the normal components having
$S=\varnothing$.

Set $C_\bullet(p)=\mathcal C_\bullet/\varepsilon\mathcal C_\bullet$.
The splitting and the contracting homotopies survive specialization.
Each degree-$k$ term has dimension $\binom Lk\nu_k(p)$, and its
degree-zero homology is the asserted $M_p$. Euler characteristic
proves~\eqref{eq:geometric-local-positivity}.
\end{proof}

\paragraph{Why normalization cannot simply be omitted.}
Let $A_k$ denote the original local algebra of $J_k$ at $Q_k$.
Factor insertion gives quotient maps $A_0\twoheadrightarrow A_k$.
Each $A_k$ is an Artinian complete intersection and hence Gorenstein,
with a one-dimensional socle; see \cite[Section~47.21]{Stacks}.
Thus no $A_0$-linear resolution with terms
$A_k^{\binom Lk}$ can be exact in positive degrees when $a\ge1$:
its top injection would require an injection of socles of dimensions
$\binom La$ and $\binom L{a-1}$, whereas $L\ge2n\ge4a$ makes
the first strictly larger. Taking the vector-space duals of these
algebras does not remove this obstruction. The special fibers of
$\mathcal B_I$ have the same lengths as $A_k$, but can have
different module structures, so the normalized construction avoids
this problem.

For example, for $n=2$ and $Q_0=t^2-1$, there are $L$ nearby
sections $Q=D_i$. Their normalizations contribute one copy of $T$
each, and each cancels with its degree-one term. The possible
remaining component has special-fiber length
$\nu_0-L\nu_1$. This is a geometric cancellation of branches,
rather than an injection of $L$ copies of the residue field into
the original local algebra.

For $n=4$ at $Q_0=(t^2-1)^2$, the local lengths already computed
in Proposition~\ref{prop:n4-double-pair} and the subsequent reduced
calculation make the normalized complex especially explicit. There
is a section $Q=D_iD_j$ for each unordered pair of distinct sites.
For even $L$, these $\binom L2$ sections exhaust the rank $\nu_0$,
and the complex is the direct sum of $\binom L2$ blocks
\[
 0\longrightarrow T\xrightarrow{\,(-1,1)\,}T^2
       \xrightarrow{\,(1,1)\,}T\longrightarrow0.
\]
For odd $L$, the same sections account for $\binom L2$ of the
$\binom{L+1}2$ branches. For each fixed site $i$, its degree-one
family has rank $L$: $L-1$ sections contain a second site pair,
and the remaining rank-one component contains only the pair at $i$.
These $L$ components are distinct and exhaust the remaining rank
of the degree-zero family. Thus there are also $L$ blocks
$0\to T\xrightarrow{1}T\to0$ in degrees one and zero. Both
parities give an acyclic complex, although their individual term
dimensions differ.

The geometric construction proves nonnegativity and gives an explicit
homology space. It does not identify that space with a spin-chain
eigenspace, prove its dimension is at most one, or establish vanishing
at every repeated singular pair.

\subsection{A spectral interpretation of the alternating local length}
\label{sec:balanced-spectral-interpretation}

The transfer matrices give a stronger interpretation of the same local
integer, including the upper bound one on its value.
Let $\tau_p$ be the polynomial quotient in the full $TQ$ relation
for $Q_0=(t^2-1)^aH$. Write $\eta_{L,n}(\tau_p)$ for the dimension
of the joint generalized eigenspace with eigenvalue polynomial
$\tau_p$ in the highest-weight subspace of $\mathcal H$ with $n$
down spins; set this dimension to zero when the eigenvalue is absent.

\begin{proposition}\label{prop:balanced-spectral-euler}
Under the standing finiteness assumptions,
\begin{equation}\label{eq:balanced-spectral-euler}
 \sum_{k=0}^a(-1)^k\binom Lk\nu_k(p)
 =\eta_{L,n}(\tau_p)\in\{0,1\}.
\end{equation}
Thus the desired cancellation is equivalent to excluding $\tau_p$
from the spectrum of the finite spin chain.
The same identity holds for $a=0,1$ and when $H$ has repeated roots
away from $\pm1$, provided $Q_0$ satisfies the full $TQ$ relation.
No root-simplicity assumption is used in this proposition.
\end{proposition}

\begin{proof}
Let $M_s$ be the highest-weight $\mathfrak{sl}_2$ Verma module of
spin $s$, with generators $e,f,h$ normalized by $[e,f]=2h$.
The singular vector $f^2v_{1/2}$ gives the exact sequence
\[
 0\longrightarrow M_{-3/2}\longrightarrow M_{1/2}
   \longrightarrow\C^2\longrightarrow0.
\]
Equip each factor with the evaluation Lax matrix whose highest-vector
diagonal entries are $t+2s$ and $t-2s$. The sequence intertwines these
matrices. Tensoring $L$ copies therefore gives an exact complex
compatible with the transfer matrices. Its term in degree $k$ is
the direct sum over $k$-element subsets of sites, with $M_{-3/2}$ at
these sites and $M_{1/2}$ at the others.

Fix $h_0=L/2-n\ge0$. In degree $k$, this weight is reached after
$n-2k$ lowering operations. Passing to highest vectors of weight
$h_0$ preserves exactness. Indeed, on this weight space the finite
sum
\[
 \mathcal P_{h_0}
 =\sum_{j\ge0}\frac{(-1)^j f^je^j}
                 {j!\prod_{q=1}^j(2h_0+q+1)}
\]
is a natural idempotent onto $\ker e$. Its denominators are nonzero,
and it commutes with the differentials and transfer matrices.
Taking its image is consequently taking a direct summand of the
exact weight-space complex.

The vacuum polynomials in degree $k$ are
\[
 a_k(t)=(t+1)^{L-k}(t-3)^k,\qquad
 d_k(t)=(t-1)^{L-k}(t+3)^k.
\]
The spectral multiplicities on each such highest-weight space equal
the local lengths of $J_k$. Here is the specialization argument
needed for this statement. Keep the local spins fixed and vary the
evaluation parameters $\theta_j$. For widely separated parameters,
suitably rescaled creation operators $\mathcal B(\theta_j)$ have leading
terms $f_j$. Their products give a basis in each fixed finite range
of degrees, by comparison with the tensor-product Verma basis.
The cyclic map from the holomorphic Yangian module to the tensor
product is therefore an isomorphism in these degrees for generic
parameters; its intertwining property follows from the defining
highest-vector and creation-operator relations in
\cite[Proposition~2.5.1]{MTV2009}. On its highest-weight part, the
transfer action is the dual regular representation of the
polynomial-solution algebra \cite[Corollary~7.3.2]{MTV2009}.
This algebra is the remainder algebra $J_k$, after translating and
rescaling the spectral variable. Traces of polynomials in the transfer
coefficients consequently agree generically. The highest-weight
projectors are regular near the homogeneous parameters. Finite
flatness of these remainder families was proved directly in
Proposition~\ref{prop:geometric-remainder-flatness}; compare also
\cite[Section~4.2]{MTV2009}. The trace identities specialize
to these parameters. Hence the joint spectral multiplicities agree
even if the specialized modules themselves are not isomorphic.
The separating condition here is $L-2n+1+j\ne0$ for the relevant
positive integers $j$, which follows from $L\ge2n$.

All roots of $J_k$ with eigenvalue $\tau_p$ must be $Q_k$. In fact,
if a monic polynomial $Q$ of degree $n-2k$ has that reduced
eigenvalue, then $(t^2-1)^kQ$ solves the full $TQ$ relation with
eigenvalue $\tau_p$. The uniqueness argument in the proof of
Proposition~\ref{prop:one-pair-weight} applies without any assumption
on root multiplicities, so this polynomial must equal $Q_0$.
Thus the multiplicity in degree $k$ is $\nu_k(p)$, and it is zero
for $k>a$.

Taking a joint generalized eigenspace is an exact operation on these
finite-dimensional modules: it is a primary direct summand for the
commuting transfer coefficients. Euler characteristic of the resulting
exact complex gives~\eqref{eq:balanced-spectral-euler}, including
the factor $\binom Lk$ counting subsets of sites. Finally, the
homogeneous finite spin chain has simple joint spectrum on highest
vectors \cite[Theorem~12.1.1]{MTV2009}, so its multiplicity is zero
or one.
\end{proof}

\subsection{Wronskian residue obstructions}
\label{sec:balanced-wronskian-obstructions}

The polynomial-kernel characterization of the finite-chain spectrum
\cite[Section~12.8]{MTV2009} identifies
$\eta_{L,n}(\tau_p)=1$ with the existence of a polynomial $P$ such
that
\begin{equation}\label{eq:balanced-polynomial-companion}
 P(t+1)Q_0(t-1)-P(t-1)Q_0(t+1)=t^L.
\end{equation}
This is stronger than the single $TQ$ relation. The following gives
explicit tests for failure of this condition.

Set
\[
 \mathcal U(t)=\frac{t^L}{Q_0(t+1)Q_0(t-1)}
 =\frac{t^{L-2a}}{(t^2-4)^aH(t+1)H(t-1)}.
\]
For $1\le j\le a$, let $c_j^+$ and $c_j^-$ be the coefficients of
$(t-2)^{-j}$ and $(t+2)^{-j}$ in its principal parts at $2$ and
$-2$, respectively, and put $\Theta_j=c_j^++c_j^-$.

\begin{proposition}\label{prop:balanced-residue-obstructions}
A polynomial companion in~\eqref{eq:balanced-polynomial-companion}
exists if and only if
\[
 \Theta_1=\cdots=\Theta_a=0.
\]
Consequently, for $a\ge2$, any nonzero $\Theta_j$ proves the desired
vanishing of the corrected weight at $p$.
\end{proposition}

\begin{proof}
Writing $f=P/Q_0$, the Wronskian relation becomes
$f(t+1)-f(t-1)=\mathcal U(t)$. Its rational principal parts telescope
along each orbit of translation by $2$. At the ordinary roots $z_j$,
the two simple-pole coefficients of $\mathcal U$ at $z_j+1$ and
$z_j-1$ sum to zero by the $TQ$ equation at $z_j$.
These contributions are the differences of rational functions with
simple poles at $z_j$.

The remaining poles are those at $2$ and $-2$, of orders at most $a$.
There is no pole at zero since $L\ge2n$. For each order $j$, their
sum of coefficients must vanish. Conversely, when it vanishes, the
corresponding principal parts are the difference of
\[
 -c_j^+\bigl((t-1)^{-j}+(t+1)^{-j}\bigr).
\]
Every polynomial part also has a polynomial discrete antiderivative.
We therefore obtain $f$ with poles only at roots of $Q_0$, and of
orders bounded by their multiplicities. Hence $P=Q_0f$ is a
polynomial. This proves the equivalence. The last assertion follows
from Proposition~\ref{prop:balanced-spectral-euler} and
\eqref{eq:balanced-local-normalization}.
\end{proof}

The highest principal-part coefficient is
\[
 c_a^+=\frac{2^{L-4a}}{H(3)H(1)},\qquad
 c_a^-=\frac{(-1)^{L-a}2^{L-4a}}{H(-1)H(-3)}.
\]
Thus $\Theta_a=0$ is exactly $\Xi(p)=0$. When this holds, write
$\mathfrak h(t)=H'(t)/H(t)$ and define
\begin{equation}\label{eq:balanced-second-obstruction}
 K(p)=\frac{2L-5a}{2}-\mathfrak h(3)-\mathfrak h(1)
                         +\mathfrak h(-1)+\mathfrak h(-3).
\end{equation}
Differentiating the pole numerators gives
\[
 \Theta_{a-1}=c_a^+K(p).
\]
Therefore $\Xi(p)=0$ and $K(p)\ne0$ also prove cancellation, with
no simplicity or Jacobian assumption on the remaining roots.

\subsection{Uniform restrictions on a possible counterexample}

The second residue obstruction also constrains configurations with
arbitrarily many additional roots. Count these roots with their
multiplicities and write
\[
 Q(t)=(t^2-1)^a\prod_{j=1}^r(t-z_j),\qquad
 z_j\ne\pm1,\qquad C=\frac{2L-5a}{2}>0.
\]

\begin{proposition}\label{prop:balanced-counterexample-localization}
Let $a\ge2$, $L\ge2n$, and suppose the corrected coefficient of
this balanced configuration is nonzero. Then
\begin{equation}\label{eq:balanced-counterexample-sum}
 C=\sum_{j=1}^r\left(\frac2{1-z_j^2}+\frac6{9-z_j^2}\right).
\end{equation}
In particular $r>0$, and at least one additional root satisfies
\begin{equation}\label{eq:balanced-counterexample-distance}
 \min\{|z_j^2-1|,|z_j^2-9|\}
 \le\frac{8r}{C}=\frac{16r}{2L-5a}.
\end{equation}
If all the additional roots are purely imaginary, a necessary
condition is
\begin{equation}\label{eq:balanced-real-rapidity-bound}
 L\le\frac52a+\frac83r.
\end{equation}
These conclusions do not require the additional roots to be distinct.
\end{proposition}

\begin{proof}
By Proposition~\ref{prop:exact-corrected-weights}, a nonzero
coefficient gives a polynomial companion. The Wronskian identity
forces $P,Q$ to be coprime and $Q(\pm3)\ne0$: a common root would
give two distinct zeros of $t^L$, while evaluation at $\pm2$
excludes the latter values when $Q(\pm1)=0$.

The pole comparison for $P/Q$ at $0,2,-2$ gives
$\Theta_a=\Theta_{a-1}=0$, hence $K=0$. This comparison is local
at the singular pair and is unchanged by multiplicities elsewhere.
Since
\[
 \frac{H'(t)}{H(t)}=\sum_{j=1}^r\frac1{t-z_j},
\]
substitution into~\eqref{eq:balanced-second-obstruction} gives
\eqref{eq:balanced-counterexample-sum}. For $r=0$ it contradicts
$C>0$. If both distances in
\eqref{eq:balanced-counterexample-distance} were greater than
$8r/C$ for every root, the triangle inequality would make the
absolute value of the sum in~\eqref{eq:balanced-counterexample-sum}
strictly less than $r(2+6)C/(8r)=C$, again a contradiction.

Finally, if $z_j=\ii y_j$ with $y_j\in\mathbb R$, each summand is
positive and at most
\[
 \frac2{1+y_j^2}+\frac6{9+y_j^2}\le\frac83.
\]
Thus $C\le8r/3$, which is
\eqref{eq:balanced-real-rapidity-bound}.
\end{proof}

Consequently, if all remaining rapidities are real in the original
variables, the corrected weight is zero whenever
$L>5a/2+8r/3$. Under $L\ge2n=4a+2r$, this includes every
$r<9a/4$. This is an exclusion criterion for that class, not an
assumption that all physical roots are real.

There is also a useful reality constraint on every possible
counterexample:
\[
 Q(t)=(-1)^n\overline{Q(-\bar t)},\qquad
 H(t)=(-1)^r\overline{H(-\bar t)}.
\]
Indeed, undoing the rescaling gives a polynomial Wronskian plane
with imaginary step $\ii$ and normalized Wronskian $\lambda^L$.
The reality theorem of \cite[Theorem~1.1]{MTV2011} makes this plane
real. Its unique monic minimal-degree polynomial therefore has real
coefficients in $\lambda$, giving the displayed identities. Thus
nonimaginary ordinary roots in the scaled variables occur in pairs
$z,-\bar z$; they are not excluded by the real-rapidity criterion.

\begin{proposition}[A bound from all residue obstructions]
\label{prop:balanced-multiplicity-bound}
For a balanced configuration with $a\ge2$ and $r=n-2a$ additional
roots, a nonzero corrected coefficient requires
\begin{equation}\label{eq:balanced-multiplicity-bound}
 a\le4r+3.
\end{equation}
At half filling $L=2n$, the stronger necessary condition is
\begin{equation}\label{eq:balanced-half-filling-bound}
 a\le4r+1.
\end{equation}
In particular, violation of either applicable bound proves vanishing,
with no reality or simplicity condition on the additional roots.
\end{proposition}

\begin{proof}
A nonzero coefficient gives a polynomial companion, so all the
singular residue obstructions vanish. Put $M=L-2a$ and introduce
the two units at $s=0$
\[
 A(s)=\frac{(s+2)^M}{(s+4)^aH(s+3)H(s+1)},\qquad
 B(s)=\frac{(s-2)^M}{(s-4)^aH(s-1)H(s-3)}.
\]
They are the analytic pole numerators of $\mathcal U$ at $2$ and
$-2$. The equations $\Theta_1=\cdots=\Theta_a=0$ say exactly
$A+B=O(s^a)$. Thus the rational logarithmic derivative
\[
 \begin{aligned}
 D(s)&=\frac{(A/B)'(s)}{(A/B)(s)}\\
 &=-\frac{4M}{s^2-4}+\frac{8a}{s^2-16}
   +\frac{H'(s-1)}{H(s-1)}+\frac{H'(s-3)}{H(s-3)}
   -\frac{H'(s+1)}{H(s+1)}-\frac{H'(s+3)}{H(s+3)}
 \end{aligned}
\]
has a zero of order at least $a-1$ at zero.

This rational function is not identically zero. Since
$H(\pm1)H(\pm3)\ne0$, its residue at $s=2$ is
\[
 -M-\operatorname{mult}_{5}(H)<0.
\]
Here multiplicities are nonnegative and $M>0$. A common denominator is
\[
 T(s)=(s^2-4)(s^2-16)
       H(s-1)H(s-3)H(s+1)H(s+3),
\]
of degree $4r+4$, with $T(0)\ne0$. Also $D(s)=O(s^{-2})$ at
infinity, so the nonzero polynomial $N=T D$ has degree at most
$4r+2$. Its order at zero is at least $a-1$, proving
\eqref{eq:balanced-multiplicity-bound}.

For the improvement at half filling, expansion at infinity gives
\[
 D(s)=-\frac{4(L-2n)}{s^2}
       +\frac{16\sum_{j=1}^r z_j}{s^3}+O(s^{-4}).
\]
When $L=2n$, comparison of the next-to-leading Wronskian coefficient
gives $[t^{n-1}]Q=0$, hence $\sum_j z_j=0$.
Consequently $D(s)=O(s^{-4})$ and $\deg N\le4r$.
The same order comparison yields $a-1\le4r$, proving
\eqref{eq:balanced-half-filling-bound}.
\end{proof}

\subsection{Two additional roots: $n=2a+2$}

\begin{proposition}\label{prop:balanced-two-additional-roots}
Let $a\ge2$, $n=2a+2$, and $L\ge4a+4$. The corrected weight is zero
at every balanced configuration with two distinct additional roots
outside $\pm1$.
\end{proposition}

\begin{proof}
By the preceding propositions, it suffices to exclude
$\Xi(p)=K(p)=0$. Write $H(t)=t^2-st+p$ and put
$C=(2L-5a)/2$. All four values $H(\pm1),H(\pm3)$ are nonzero.
The resonance condition is
\begin{equation}\label{eq:two-additional-resonance}
 \begin{cases}
 (p+9)(p+1)+3s^2=0,&L\equiv a\pmod2,\\
 s(p+3)=0,&L\not\equiv a\pmod2.
 \end{cases}
\end{equation}
For later use, direct evaluation of logarithmic derivatives gives
\[
 K=C-\frac{6(2p+18-s^2)}{(p+9)^2-9s^2}
       -\frac{2(2p+2-s^2)}{(p+1)^2-s^2}.
\]

The product of the reduced Bethe equations, combined with resonance,
also implies
\begin{equation}\label{eq:balanced-resonant-momentum}
 \Pi^L=(-1)^{aL},\qquad
 \Pi=\frac{H(-1)}{H(1)}=\frac{p+1+s}{p+1-s}.
\end{equation}
Indeed, with $\Sigma=H(-3)/H(3)$ these two conditions are
$\Pi^{L-a}=\Sigma^a$ and $\Pi\Sigma=(-1)^{a+L+1}$.
If $p,s$ are real, \eqref{eq:balanced-resonant-momentum} forces
$\Pi=\pm1$, hence $s=0$ or $p=-1$.

If $L\equiv a\pmod2$, substitute
$s^2=-(p+9)(p+1)/3$ into $K=0$. This gives
\[
 p=-3+\frac8{C-1},\qquad
 s^2=-\frac{(p+9)(p+1)}3.
\]
Here $C\ge(3a+8)/2\ge7$, so $-3<p<-1$ and $s^2>0$.
Thus $p,s$ are real with $s\ne0$ and $p\ne-1$, contradicting
\eqref{eq:balanced-resonant-momentum}.

Now let $L\not\equiv a\pmod2$. If $p=-3$, the equation $K=0$
gives
\[
 s^2=\frac{12C}{3C-8}>0.
\]
Again $p,s$ are real, $s\ne0$, and $p\ne-1$, giving the same
contradiction. It remains to consider $s=0$.

For $H=t^2+p$, the equation $K=0$ becomes
\begin{equation}\label{eq:two-additional-quadratic-obstruction}
 F(p):=C(p+1)(p+9)-16p-48=0.
\end{equation}
This quadratic is irreducible over $\mathbb Q$. Its discriminant is
$16(d^2+12)$, where $d=2C-2=2L-5a-2$ is an integer at least $12$.
Since $d^2<d^2+12<(d+1)^2$, the discriminant is not a rational
square. The two real conjugate roots of $F$ lie respectively in
$(-1,0)$ and $(-9,-3)$: this follows from
\[
 F(-1)=-32,\quad F(0)=9C-48>0,\quad
 F(-9)=96,\quad F(-3)=-12C.
\]

The conditions $H\mid S_a$ are polynomial equations in $p$ with
rational coefficients. If either root of the irreducible $F$ satisfied
them, both would. Choose the conjugate $p\in(-1,0)$ and set
$w=\sqrt{-p}\in(0,1)$, so the additional roots are $w,-w$.
The reduced equation at $w$, after cancelling its interaction with
$-w$, requires
\[
 \left(\frac{w+1}{w-1}\right)^{L-a-1}
 =\left(\frac{w+3}{w-3}\right)^a.
\]
But
\[
 \left|\frac{w+1}{w-1}\right|
 >\left|\frac{w+3}{w-3}\right|>1,
 \qquad L-a-1>a,
\]
so equality is impossible. This excludes the final case and proves
the proposition.
\end{proof}

\subsection{Three additional roots invariant under sign reversal}

\begin{proposition}\label{prop:balanced-three-symmetric-roots}
Let $a\ge2$, $n=2a+3$, and $L\ge4a+6$. The corrected weight
vanishes at every balanced configuration whose three additional roots
are $0,w,-w$, with $w\ne0$.
\end{proposition}

\begin{proof}
Write $H(t)=t(t^2+p)$, where $p=-w^2$. The values
$H(\pm1),H(\pm3)$ are nonzero, as in the preceding arguments.
Since $H$ is odd, $\Xi=0$ requires $L\not\equiv a\pmod2$.
Otherwise the highest residue obstruction already proves vanishing.

Put
\[
 C=\frac{2L-5a}{2},\qquad D=C-\frac83.
\]
The root at zero contributes $8/3$ to the logarithmic-derivative
sum in~\eqref{eq:balanced-second-obstruction}. Consequently
\[
 K=D-\frac{12}{p+9}-\frac4{p+1}.
\]
Thus $K=0$ would imply
\[
 F(p):=D(p+1)(p+9)-16p-48=0.
\]
This quadratic is irreducible over $\mathbb Q$. Indeed, its
discriminant is
\[
 16\bigl((2D-2)^2+12\bigr)
 =\frac{16}{9}(N^2+108),\qquad N=6L-15a-22\ge32.
\]
If $N^2+108$ were an integer square, factoring the difference of
the two squares would give $N=6$ or $N=26$: these are the only
nonnegative possibilities from factor pairs of $108$ with the same
parity. Both contradict $N\ge32$.

Here $D\ge19/3$, and the signs of $F$ at $-9,-3,-1,0$
show that its roots lie in $(-9,-3)$ and $(-1,0)$.
The reduced equations $H\mid S_a$ have rational coefficients in
$p$. Hence if one root of this irreducible quadratic satisfied them,
so would its conjugate $p\in(-1,0)$. Choose this conjugate and put
$w=\sqrt{-p}\in(0,1)$. The reduced equation at $w$, including its
interactions with $0$ and $-w$, becomes
\[
 \left(\frac{w+1}{w-1}\right)^{L-a-1}
 =\left(\frac{w+3}{w-3}\right)^a\frac{w+2}{w-2}.
\]
But
\[
 \left|\frac{w+1}{w-1}\right|
 >\left|\frac{w+2}{w-2}\right|
 >\left|\frac{w+3}{w-3}\right|>1,
 \qquad L-a-1>a+1,
\]
so this equality is impossible. Thus $\Xi$ and $K$ cannot both
vanish, and the residue obstruction proves the proposition.
\end{proof}

\subsection{Three additional roots: momentum and exact elimination}

For three additional roots, write
\[
 H(t)=t^3-st^2+pt-q,\qquad
 C=\frac{2L-5a}{2}\ge9,\qquad
 \Pi=\frac{H(-1)}{H(1)}.
\]
The product of the reduced equations and $\Xi=0$ again imply
$\Pi^L=(-1)^{aL}$. The resonance equation itself is
\begin{equation}\label{eq:cubic-resonance}
 \begin{cases}
 3(p+9)(p+1)+(9s+q)(s+q)=0,&L\equiv a\pmod2,\\
 3s(p+3)+q(p+7)=0,&L\not\equiv a\pmod2.
 \end{cases}
\end{equation}

\begin{proposition}\label{prop:balanced-cubic-endpoint-momentum}
For $a\ge2$ and $L\ge4a+6$, a balanced configuration with three
additional roots and $\Pi=1$ or $\Pi=-1$ has zero corrected weight.
\end{proposition}

\begin{proof}
It suffices to consider $\Xi=0$. If $\Pi=1$, then $p=-1$ and
$s+q\ne0$, since $H(1)\ne0$. For opposite parities,
\eqref{eq:cubic-resonance} would force $s+q=0$, a contradiction.
For equal parities it gives $q=-9s$. Substitution into
\eqref{eq:balanced-second-obstruction} yields $K=C-5/3\ne0$.

If $\Pi=-1$, then $q=-s$ and $p\ne-1$. For equal parities,
\eqref{eq:cubic-resonance} gives $p=-9$, and direct substitution
gives $K=C-3\ne0$. For opposite parities, the resonance equation
gives $s=0$, hence $q=0$. This is the sign-invariant case of
Proposition~\ref{prop:balanced-three-symmetric-roots}.
\end{proof}

For the remaining momentum values, put
\[
 v=\frac{s+q}{p+1},\qquad
 \Pi=-\frac{1+v}{1-v}.
\]
The endpoint proposition allows us to assume $v\ne0$ and $p\ne-1$.
Equations~\eqref{eq:cubic-resonance} then give
\begin{equation}\label{eq:cubic-momentum-parameterization}
 s=\begin{cases}
 -\dfrac{3(p+9)+v^2(p+1)}{8v},&L\equiv a\pmod2,\\[5pt]
 -\dfrac{v(p+7)}2,&L\not\equiv a\pmod2,
 \end{cases}
 \qquad q=v(p+1)-s.
\end{equation}
The momentum condition is the one-variable polynomial equation
\[
 (-(1+v))^L-(-1)^{aL}(1-v)^L=0.
\]
Over each number field defined by an irreducible factor of this
polynomial, $K=0$ is a polynomial equation of degree at most two in
$p$. Reduce the three coefficients of $S_a\bmod H$ modulo this
quadratic. Their common gcd with the quadratic tests whether the
reduced Bethe equations can also hold. Factors supported on
$H(\pm1)H(\pm3)\operatorname{disc}(H)=0$ must be removed, since
these are outside the configurations under consideration.

\paragraph{An exact check beyond sign-invariant configurations.}
For $a=2$, $L=14$, the two nonendpoint momentum factors are
\[
 v^6+21v^4+35v^2+7,\qquad
 7v^6+35v^4+21v^2+1.
\]
In either field, the numerator of $K$, up to a nonzero scalar, is
\[
 (7v^2-3)p^2+(58v^2-18)p+51v^2+81.
\]
The gcd calculation just described gives $1$ in both fields.
Together with the endpoint proposition, this excludes every
three-additional-root configuration for these parameters. The
calculation is implemented with exact rational and number-field
arithmetic in \texttt{computations/check\_balanced\_cubic.py}; it
also verifies all lengths $14\le L\le25$ for $a=2$,
$18\le L\le25$ for $a=3$, and $22\le L\le25$ for $a=4$.
These finite checks do not establish the assertion for arbitrary $L$.

\subsection{Modular certificates for repeated singular pairs}

The companion condition also permits exact tests which do not assume
that the additional roots are distinct. The finiteness statement below
is essential: an empty modular fiber of an arbitrary affine system
does not by itself exclude complex solutions.

\begin{proposition}\label{prop:finite-wronskian-family}
Fix $L\ge2n$ and put $d=L-n+1$ and
$R=\mathbb Z[1/(2(L+1)!)]$. Consider the coefficient scheme of
\[
 P(t+1)Q(t-1)-P(t-1)Q(t+1)=t^L,
\]
where $Q$ is monic of degree $n$, $P$ has degree $d$ and leading
coefficient $1/(2(d-n))$, and $[t^n]P=0$. This scheme is finite
over $R$. Consequently, for any prime $\mathfrak p>L+1$, if this
system together with $(t^2-1)^2\mid Q$ has no solution over
$\overline{\mathbb F}_{\mathfrak p}$, it has no complex solution.
\end{proposition}

\begin{proof}
The leading coefficient of $P$ is forced by the Wronskian equation.
Adding a multiple of $Q$ to $P$ does not change that equation, so
$[t^n]P=0$ fixes this freedom uniquely. Give the coefficients of
$t^j$ in $Q$ and $P$ weights $n-j$ and $d-j$, respectively.
The highest-weight parts of the remaining coefficient equations
are those of the continuous Wronskian identity
\[
 2(P'Q-PQ')=t^L.
\]
Indeed, the first-order terms in the two shifts give this expression;
every higher odd-order term has strictly smaller coefficient weight.

Over any field of characteristic zero or greater than $L+1$, the
continuous system has only the solution
\[
 Q=t^n,\qquad P=\frac{t^d}{2(d-n)}.
\]
To see this, its gcd is $t^g$ up to a scalar, since its square divides
$t^L$. Write $P=t^g A$, $Q=t^g B$, with $A,B$ coprime.
Choose either $A/B$ or $B/A$ to have denominator nonzero at zero.
Its derivative has order $L-2g$ there. The first nonconstant term
therefore has order $L-2g+1$, while its numerator after subtracting
its value at zero has degree at most $d-g$. These orders behave as
in characteristic zero because all relevant polynomial degrees are
smaller than the characteristic. Thus
\[
 L-2g+1\le d-g,\qquad g\ge n.
\]
Since $g\le n$, we have $Q=t^n$; integrating the continuous identity
and imposing the gauge gives the displayed $P$.

Thus the homogeneous initial equations have only the coefficient
origin as their common zero in every residue characteristic of $R$.
Their radical contains every coefficient variable, so their quotient
is finite over $R$. The associated graded algebra of the discrete
coefficient scheme is a quotient of this finite algebra. The
discrete coefficient algebra is consequently finite over $R$ as well.

Imposing $(t^2-1)^2\mid Q$ gives a further finite $R$-algebra $B$.
If $B/\mathfrak pB=0$, Nakayama's lemma gives
$B\otimes_R R_{(\mathfrak p)}=0$, and hence
$B\otimes_R\mathbb Q=0$. This proves the modular implication.
\end{proof}

\paragraph{Eliminating the companion coefficients.}
This test can be performed using only the coefficients of
$Q=(t^2-1)^2H$. Descend through the coefficients of $P$, starting
with degree $d$. For $m\ne n$, the coefficient of $t^{m+n-1}$ in
\[
 (t+1)^mQ(t-1)-(t-1)^mQ(t+1)
\]
is the nonzero constant $2(m-n)$. It therefore determines
$[t^m]P$ from the coefficients already found. At $m=n$, set the
gauge coefficient to zero and retain the resulting constraint.
The remaining Wronskian coefficients give polynomial equations in
$H$. When $L=2n$, the first constraint also forces
$[t^{n-1}]Q=0$, reducing the number of variables by one.

The exact implementation
\texttt{computations/check\_balanced\_wronskian.py} obtains the unit
Gr\"obner basis modulo $101$ for
\[
 (L,n)=(16,8),\quad(17,8),\quad(18,9).
\]
Proposition~\ref{prop:finite-wronskian-family} therefore excludes a
polynomial companion for every $Q$ divisible by $(t^2-1)^2$ in
these cases. This includes all $a\ge2$ allowed by each $n$, without
an assumption on the other roots. By
Proposition~\ref{prop:balanced-spectral-euler}, every balanced
configuration under consideration has zero corrected weight for
these parameter pairs. These certificates are finite checks, not a
uniform proof for arbitrary $L,n$.

\subsection{The corrected sum without root-simplicity assumptions}

The spectral calculation gives an exact interpretation of the entire
corrected formal sum, including configurations with arbitrary
collisions. This interpretation does not require the distinct-root
assumption from Section~\ref{sec:general-tq}.
For an ordered configuration $p$, write
\[
 \sigma(p)=|\operatorname{Stab}_{S_n}(p)|
          =\prod_{z\in\C}m_z(p)!,\qquad
 Q_p(t)=\prod_{i=1}^n(t-p_i),
\]
where $m_z(p)$ is the multiplicity of the value $z$ in $p$.
Let $\varepsilon(Q)$ be one if $Q$ has a polynomial companion in
\eqref{eq:balanced-polynomial-companion}, and zero otherwise.

\begin{proposition}[Exact corrected weights]
\label{prop:exact-corrected-weights}
Assume $L\ge2n$ and the standing finiteness hypotheses on the
partition ideals. For every $p\in X$,
\begin{equation}\label{eq:exact-corrected-weight}
 \boxed{\operatorname{coeff}_p Z^{\mathrm{cand}}_{L,n}
       =\frac{\sigma(p)}{n!}\,\varepsilon(Q_p).}
\end{equation}
In particular, all corrected coefficients are nonnegative, even
without assuming simplicity of the ordinary roots.
\end{proposition}

\begin{proof}
If $p$ belongs to the support of some $Z_k^{\mathrm{dep}}$, the
depermutation correspondence implies the reduced $TQ$ relation for
the polynomial remaining after removing its $k$ fixed pairs.
Multiplication by $(t^2-1)^k$ lifts this relation to the full $TQ$
relation for $Q_p$. Hence, if $Q_p$ does not satisfy $TQ$, all the
corrected coefficients at $p$ vanish; a companion is also impossible.

Suppose now that $Q_p$ satisfies $TQ$. Comparison of orders at
$1,-1$ forces equal singular multiplicities, say $a=b$, with $a=0$
allowed. Write
\[
 Q_p=(t^2-1)^a\prod_{j=1}^m(t-z_j)^{s_j},\qquad
 z_j\ne\pm1,
\]
where the $z_j$ are distinct, but the $s_j$ are arbitrary.
For $0\le k\le a$, let $\nu_k$ be the local reduced remainder
length at $Q_k=Q_p/(t^2-1)^k$. The same root-map calculation as in
\eqref{eq:balanced-local-normalization} gives
\[
 \begin{aligned}
 c_k(p)
 &=\frac{(k!)^2(n-2k)!}{n!}\binom ak^2
   \frac{((a-k)!)^2\prod_j s_j!}{(n-2k)!}\,\nu_k\\
 &=\frac{(a!)^2\prod_j s_j!}{n!}\,\nu_k
  =\frac{\sigma(p)}{n!}\,\nu_k.
 \end{aligned}
\]
For $k>a$ the coefficient is zero. The proof of
Proposition~\ref{prop:balanced-spectral-euler} applies verbatim:
neither the specialization of the Verma-module characters nor the
uniqueness of the fixed-eigenvalue polynomial uses separability of
$H$. Thus the alternating sum of the $\nu_k$ is
$\eta_{L,n}(\tau_p)$. The polynomial-kernel characterization of the
finite-chain spectrum identifies this number with
$\varepsilon(Q_p)$, proving the formula.
\end{proof}

Let $\mathcal Q_{L,n}$ be the finite set of monic degree-$n$
polynomials which admit a normalized polynomial companion, and let
$J^{\mathrm{poly}}_{L,n}\subset\C[e_1,\ldots,e_n]$ be its reduced
coefficient ideal. For the ordered-root map
$q:\A^n\to\Sym^n(\A^1)$, formula~\eqref{eq:exact-corrected-weight}
is equivalently the formal-sum identity
\begin{equation}\label{eq:corrected-polynomial-kernel-ideal}
 Z^{\mathrm{cand}}_{L,n}
 =\frac1{n!}Z\bigl(q^*J^{\mathrm{poly}}_{L,n}\bigr),\qquad
 q_*Z^{\mathrm{cand}}_{L,n}=\sum_{Q\in\mathcal Q_{L,n}}[Q].
\end{equation}
Indeed, the local degree of the root map at $p$ is $\sigma(p)$.
An orbit has $n!/\sigma(p)$ ordered representatives, so its total
corrected weight is exactly $\varepsilon(Q_p)$.

\begin{proposition}[Affine completeness degree]
\label{prop:affine-completeness-degree}
Under the finiteness hypotheses, without the distinct-root assumption,
\[
 \deg Z^{\mathrm{cand}}_{L,n}
 =\binom Ln-\binom L{n-1}.
\]
\end{proposition}

\begin{proof}
The depermutation correspondence gives
$\deg Z_k^{\mathrm{dep}}=\ell(J_k)$: passing to the ordered roots
and applying the prescribed normalization preserves the total weight.
Proposition~\ref{prop:geometric-remainder-flatness}, specialized at
$\theta=0$, gives
\[
 \ell(J_k)=\binom{L+n-2k-2}{n-2k}.
\]
Summing with coefficients $(-1)^k\binom Lk$ and applying
\eqref{eq:virtual-completeness-degree} proves the assertion. This
degree argument does not use a spectral theorem.
\end{proof}

Thus the affine degree assertion in
Conjecture~\ref{conj:alternating-completeness} follows from the
geometric argument above. The pointwise assertion requires more: a
polynomial-kernel solution must have the root-simplicity properties
prescribed in the target weights. The degree identity alone cannot
exclude a solution with repeated singular pairs.

\subsection{The uniform statement needed to finish the proof}

Let $\mathcal W^{(2)}_{L,n}$ be the closed subspace of the normalized
Wronskian coefficient scheme in
Proposition~\ref{prop:finite-wronskian-family} defined by
\[
 (t^2-1)^2\mid Q(t).
\]
There is no restriction on the other factors of $Q$. In particular,
this single divisibility condition includes every $a\ge2$.

\begin{proposition}[Uniform reduction to double-pair exclusion]
\label{prop:balanced-uniform-reduction}
Fix $L\ge2n$, $n\ge4$, and retain only the finiteness hypotheses
of this note. The following statements are equivalent:
\begin{enumerate}
 \item Every balanced configuration with $a=b\ge2$ has zero
 corrected weight in $Z^{\mathrm{cand}}_{L,n}$.
 \item $\mathcal W^{(2)}_{L,n}(\C)=\varnothing$.
\end{enumerate}
\end{proposition}

\begin{proof}
By Proposition~\ref{prop:exact-corrected-weights}, the corrected
coefficient is $\sigma(p)\varepsilon(Q_p)/n!$. Hence the second
statement implies the first, regardless of collisions among the
ordinary roots.

Conversely, let $(P,Q)$ be a complex point of
$\mathcal W^{(2)}_{L,n}$. Two elementary consequences of its
Wronskian equation identify its roots with the configurations above.
First, $P,Q$ are coprime: a common root $c$ would make the
Wronskian vanish at both $c+1$ and $c-1$, forcing $c=-1$ and
$c=1$ simultaneously. Consequently $P(1),P(-1)\ne0$.
If the multiplicities of $1,-1$ in $Q$ were unequal, the
Wronskian at zero would have the smaller of these orders, which is
less than $L$. Hence these multiplicities are equal, say $a\ge2$.

Second, using the Wronskian at $t+1$ and $t-1$ gives
\[
 \begin{aligned}
 &(t+1)^LQ(t-2)+(t-1)^LQ(t+2)\\
 &\quad=Q(t)\bigl(P(t+2)Q(t-2)-P(t-2)Q(t+2)\bigr).
 \end{aligned}
\]
Thus $Q$ satisfies the full polynomial $TQ$ relation, so labeling
its roots gives a balanced point $p\in X$. Its companion implies
$\varepsilon(Q_p)=1$, and its corrected coefficient is
$\sigma(p)/n!>0$, contradicting the first statement.
\end{proof}

The unresolved assertion is therefore precisely
\[
 \boxed{\mathcal W^{(2)}_{L,n}(\C)=\varnothing
        \quad\text{for every }L\ge2n.}
\]
Equivalently, the minimal-degree polynomial in a polynomial
Wronskian pair should have at most a simple root at each of $1,-1$.
This is a uniform assertion, rather than a separate calculation for
each number of singular pairs.

\paragraph{Why the minimal-degree restriction is essential.}
The exclusion is false for arbitrary polynomial solutions in the same
two-dimensional kernel. In fact, double singular pairs occur naturally
in the higher-degree dual pencil. The existence and role of this dual
pencil are discussed in~\cite{PS1999}; the following observation
explains a multiplicity phenomenon within it.

\begin{proposition}\label{prop:dual-double-pairs}
Let $L\ge3$ be odd. Suppose an even polynomial $Q_*$ of degree
$m\le L/2$, with $Q_*(1)\ne0$, has a polynomial companion $P_*$.
Then a member of the dual pencil $P_*+\C Q_*$ is divisible by
$(t^2-1)^2$. Its degree is $L-m+1>L/2$.
\end{proposition}

\begin{proof}
Normalize $P_*$ by $[t^m]P_*=0$. Its degree $L-m+1$ is even.
Reflection $t\mapsto-t$ preserves the normalized Wronskian equation,
since $Q_*$ is even and $L$ is odd. Uniqueness of the normalized
companion therefore makes $P_*$ even. Put
\[
 R(t)=P_*(t)-\frac{P_*(1)}{Q_*(1)}Q_*(t).
\]
Then $R$ is even and $R(1)=R(-1)=0$. Differentiating
\[
 R(t+1)Q_*(t-1)-R(t-1)Q_*(t+1)=t^L
\]
at zero gives
\[
 2\bigl(R'(1)Q_*(1)-R(1)Q_*'(1)\bigr)=0.
\]
Hence $R'(1)=0$, and evenness also gives $R'(-1)=0$.
The subtraction does not change its degree, proving the claim.
\end{proof}

For an explicit example,
\[
 L=5,\qquad Q(t)=(t^2-1)^2,\qquad
 P(t)=-\frac14(t^2+1)
\]
satisfy
\[
 P(t+1)Q(t-1)-P(t-1)Q(t+1)=t^5.
\]
The same polynomial kernel has minimal-degree solution $t^2+1$.
Thus the example has a polynomial companion and a physical transfer
eigenvalue, but the repeated-pair polynomial has degree four, beyond
the range $L\ge2\deg Q$. It is not a counterexample to the boxed
assertion. Any uniform proof must use the minimal-degree condition,
rather than polynomiality of both solutions alone.

\paragraph{What the cited results do and do not supply.}
The polynomial-kernel result of \cite[Section~12.8]{MTV2009}
identifies the Wronskian pairs with the finite-chain spectrum; it
does not by itself exclude the closed subspace
$\mathcal W^{(2)}_{L,n}$. The multiplicity constraints of
\cite[Section~4.2]{HJJ2024} imply equality of the two singular
multiplicities for spin $1/2$, but allow their common value to
exceed one. Neither statement supplies the boxed exclusion.
Even the single $TQ$ relation together with the distinct-root
assumption away from $\pm1$ is insufficient: for every $L\ge8$,
\[
 \begin{aligned}
 Q(t)&=(t^2-1)^2,\\
 \tau(t)&=(t+1)^{L-2}(t-3)^2
            +(t-1)^{L-2}(t+3)^2
 \end{aligned}
\]
satisfy $TQ$, and the condition on ordinary roots is vacuous.
The absence of a companion in these examples was established by
the pure-cluster calculation, not by $TQ$ alone.

In residue language, the remaining task is to exclude the reduced
Bethe equations together with $\Theta_1=\cdots=\Theta_a=0$.
The unproved configurations have at least three additional roots,
outside the families and parameter values already treated. A
uniform proof of the boxed assertion is still required; the finite
certificates above do not replace it.

\begin{thebibliography}{99}
\bibitem{Stacks}
The Stacks Project Authors,
\emph{The Stacks Project},
\href{https://stacks.math.columbia.edu/tag/062F}{Tag~062F},
\href{https://stacks.math.columbia.edu/tag/035E}{Tag~035E}, and
\href{https://stacks.math.columbia.edu/tag/0DW6}{Section~47.21}.

\bibitem{MTV2011}
E.~Mukhin, V.~Tarasov, and A.~Varchenko,
\emph{Reality property of discrete Wronski map with imaginary step},
Lett.\ Math.\ Phys.\ \textbf{100}, 151--160 (2012),
\href{https://arxiv.org/abs/1102.5368}{arXiv:1102.5368}.

\bibitem{PS1999}
G.~P.~Pronko and Yu.~G.~Stroganov,
\emph{Bethe equations ``on the wrong side of the equator''},
J.\ Phys.\ A: Math.\ Gen.\ \textbf{32}, 2333--2340 (1999),
\href{https://arxiv.org/abs/hep-th/9808153}{arXiv:hep-th/9808153}.

\bibitem{MTV2009}
E.~Mukhin, V.~Tarasov, and A.~Varchenko,
\emph{Bethe algebra of homogeneous XXX Heisenberg model has simple spectrum},
Commun.\ Math.\ Phys.\ \textbf{288}, 1--42 (2009),
\href{https://arxiv.org/abs/0706.0688}{arXiv:0706.0688}.

\bibitem{BK2016}
L.~Bus\'e and A.~Karasoulou,
\emph{Resultant of an equivariant polynomial system with respect to the symmetric group},
J.\ Symbolic Comput.\ \textbf{76}, 142--157 (2016),
\href{https://arxiv.org/abs/1407.2799}{arXiv:1407.2799}.

\bibitem{GJ2020}
E.~Granet and J.~L.~Jacobsen,
\emph{On zero-remainder conditions in the Bethe ansatz},
JHEP \textbf{03}, 178 (2020),
\href{https://arxiv.org/abs/1910.07797}{arXiv:1910.07797}.

\bibitem{HJJ2024}
J.~Hou, Y.~Jiang, and R.-D.~Zhu,
\emph{Spin-$s$ Rational $Q$-system},
SciPost Phys.\ \textbf{16}, 113 (2024),
\href{https://arxiv.org/abs/2303.07640}{arXiv:2303.07640}.

\bibitem{Sturmfels}
B.~Sturmfels,
\emph{Solving Systems of Polynomial Equations},
CBMS Regional Conference Series in Mathematics \textbf{97}, AMS (2002),
\href{https://web.math.berkeley.edu/~bernd/cbms.pdf}{author's manuscript}.

\bibitem{HNS2013}
W.~Hao, R.~I.~Nepomechie, and A.~J.~Sommese,
\emph{Completeness of solutions of Bethe's equations},
Phys.\ Rev.\ E \textbf{88}, 052113 (2013),
\href{https://arxiv.org/abs/1308.4645}{arXiv:1308.4645}.

\bibitem{NW2013}
R.~I.~Nepomechie and C.~Wang,
\emph{Algebraic Bethe ansatz for singular solutions},
J.\ Phys.\ A \textbf{46}, 325002 (2013),
\href{https://arxiv.org/abs/1304.7978}{arXiv:1304.7978}.

\end{thebibliography}
\end{document}
